% concatenated from hsV.tex
\documentclass[a4paper,11pt]{article} 
  %{amsart}
  
%------- Packages
\usepackage[english]{babel}
\usepackage[a4paper]{geometry}
\usepackage{amsmath}
\usepackage{amsthm}
\usepackage{amssymb}
\usepackage{bbm}
\usepackage{xcolor}
\usepackage{enumerate}
\usepackage{dsfont}
\usepackage{cancel}
\usepackage{ulem} %\sout{} to cancel text

%\usepackage[color]{showkeys}

\usepackage{hyperref}         %%%% click to contents
%\usepackage{refcheck}   %checks if labeled equations have references in the text

%%%%%%%%%%%%%%%%%%%%%% TIKZ
\usepackage{tikz}
\usetikzlibrary{fit}
\usetikzlibrary{shapes.geometric}
\usetikzlibrary{decorations.pathmorphing}

\tikzset{
point/.style={circle,fill=black,inner sep=1pt},
vertex/.style={circle,fill=black,inner sep=1.5pt},   % VERTICI
bvertex/.style={circle,fill=black,inner sep=2.8pt},
%
Bvertex/.style={circle,fill=black,inner sep=4pt}, % VERTICE GRANDE
%
specialEP/.style={rectangle,fill=white,draw,inner sep=3pt},  % RETTANGOLO
%
whitevex/.style={circle,fill=white,draw, inner sep=2pt},
linelabel/.style={sloped,above,very near start, inner sep=1pt,execute at begin node=$\scriptstyle,execute at end node=$},
baseline=(current  bounding  box.center),doubled/.style={double distance= 1pt,line width=1.5pt},
th/.style={line width=0.5 pt, gray},  %linea thin GRIGIA
med/.style={line width=1 pt}  %linea con medio spessore
}

\xdefinecolor{myred}{rgb}{0.7,0,0.2} %per definire colori personalizzati
\xdefinecolor{mypurple}{rgb}{0.6,0.2,0.4} %{0.42,0.118,0.588} 
\xdefinecolor{myblue}{rgb}{0.259,0.2,0.6}   
\xdefinecolor{myorange}{rgb}{1,0.6,0.2}   
\definecolor{orange}{rgb}{1,0.5,0}
\xdefinecolor{purpleish}{cmyk}{0.75,0.75,0,0}


%%%%%%%%%%%%%%%%%

\def\fh{{\mathfrak h}}
\def\bR{\mathbb{R}}
\def\bT{\mathbb{T}}
\def\bN{\mathbb{N}}
\def\NN{\mathbb{N}}
\def\bZ{\mathbb{Z}}
\def\cC{\mathcal{C}}
\def\cU{\mathcal{U}}
\def\cQ{\mathcal{Q}}
\def\cD{\mathcal{D}}
\def\cA{\mathcal{A}}
\def\cM{\mathcal{M}}
%\def\cP{\mathcal{P}}
\def\cW{\mathcal{W}}
\def\cV{\mathcal{V}}
\def\cO{\mathcal{O}}
\def\cF{\mathcal{F}}
\def\cG{\mathcal{G}}
\def\cL{\mathcal{L}}
\def\cJ{\mathcal{J}}
\def\cN{\mathcal{N}}
\def\cE{\mathcal{E}}
\def\cK{\mathcal{K}}
\def\cH{\mathcal{H}}
\def\cR{\mathcal{R}}
\def\cS{\mathcal{S}}
\def\cT{\mathcal{T}}
\def\eps{\varepsilon}
\def\ph{\varphi}

\def\wtph{\widetilde{\varphi_{\xi}}}
\def\wt{\widetilde}

\def\NNN{\mathbb{N}}  
\def\ZZZ{\mathbb{Z}} 
\def\RRR{\mathbb{R}} 
\def\PPP{\mathbb{P}} 
\def\EEE{\mathbb{E}} 

\def\bsA{{\boldsymbol\Phi}} \def\bsA{{\boldsymbol A}}


% --- Miscellanea
\def\Val{{\rm Val}}
\def\indic{\hbox{\raise-2pt \hbox{\indbf 1}}}

\let\dpr=\partial
\let\circa=\cong
%\let\bs=\backslash
\let\==\equiv
\let\txt=\textstyle
\let\dis=\displaystyle
\let\io=\infty
\let\0=\noindent


\def\pagina{{\vfill\eject}}
\def\*{{\hfill\break\null\hfill\break}}
\def\bra#1{{\langle#1|}}
\def\ket#1{{|#1\rangle}}
\def\media#1{{\left\langle#1\right\rangle}}
\def\bmedia#1{{\bigl\langle#1\bigr\rangle}}
\def\bmed#1#2{{\bigl\langle#1\bigr\rangle_{#2}}}
\def\ie{\hbox{\it i.e.\ }}
\def\eg{\hbox{\it e.g.\ }}



\def\unit{\mathbb{I}}  
\def\openone{{\mathbbm{1}}}
\def\braket#1#2{{\langle#1|#2\rangle}}
\def\xr#1{\xrightarrow[#1]{}}
\let\arr=\rightarrow
%
\def\allgray#1{\color[gray]{0.35}{#1}} 
\def\allblue#1{\color{blue}{#1}} 
%
% \def\red#1{\textcolor{red}{#1}} 
% \def\blue#1{\textcolor{blue}{#1}}  
% \def\purple#1{\textcolor{purple}{#1}} 

\def\PS#1{\ps^{(#1)}}
\def\TPS#1#2{\tl{\ps}^{(#1)}_{#2}}
\def\txe{\text{e}}
\def\txi{\text{in}}
\def\cst{(\text{const.})}


\def\TL#1{{\widetilde #1}}

\def\Dpr{\V\dpr\,}
\def\aps{{\it a posteriori}}
\def\lft{\left}
\def\rgt{\right}
\def\der{\hbox{\rm d}}
\def\la{{\langle}}
\def\ra{{\rangle}}
\def\norm#1{{\left|\hskip-.05em\left|#1\right|\hskip-.05em\right|}}
\def\tgl#1{\!\!\not\!#1\hskip1pt}
\def\tende#1{\,\vtop{\ialign{##\crcr\rightarrowfill\crcr
             \noalign{\kern-1pt\nointerlineskip}
             \hskip3.pt${\scriptstyle #1}$\hskip3.pt\crcr}}\,}
\def\otto{\,{\kern-1.truept\leftarrow\kern-5.truept\to\kern-1.truept}\,}
\def\fra#1#2{{#1\over#2}}

\def\sde{{\cs SDe}}
\def\wti{{\cs WTi}}
\def\osa{{\cs OSa}}
\def\ce{{\cs CE}}
\def\rg{{\cs RG}}


\def \dx{{d}x}
\def \dk{{d}k}
\def\Tr{{\rm Tr}}
\def\tr{{\rm tr}}

\def\Re{{\rm Re}\,}\def\Im{{\rm Im}\,}



\newtheorem{theorem}{Theorem}[section]  % use thm for %Theorems to keep numbering consistent
\newtheorem{corollary}[theorem]{Corollary}
\newtheorem{proposition}[theorem]{Proposition}
\newtheorem{lemma}[theorem]{Lemma}
\newtheorem{remark}{Remark}



\numberwithin{equation}{section}

%%%%% Serena's commands
\def\tl#1{{\tilde{#1}}}
\def\be{\begin{equation}}
\def\ee{\end{equation}}
\def\spl#1{\[ \begin{split}#1\end{split} \]}
\def\bes#1{\be \begin{split}#1\end{split} \be}  %numbered split
\newcommand{\hc}{\mbox{h.c.}}
\def \note#1{\marginpar{\footnotesize{#1}}}

%Greek letters
\let\a=\alpha \let\b=\beta    \let\g=\gamma     \let\d=\delta     \let\e=\varepsilon
\let\z=\zeta  \let\h=\eta     \let\th=\vartheta \let\k=\kappa     \let\l=\lambda
\let\m=\mu    \let\n=\nu      \let\x=\xi        \let\p=\pi        \let\r=\rho
\let\s=\sigma \let\t=\tau     \let\f=\phi    \let\ph=\varphi   \let\c=\chi
\let\ps=\psi  \let\y=\upsilon \let\o=\omega     \let\si=\varsigma
\let\G=\Gamma \let\D=\Delta   \let\Th=\Theta    \let\L=\Lambda    \let\X=\Xi
\let\P=\Pi    \let\Si=\Sigma  \let\F=\Phi       \let\Ps=\Psi
\let\O=\Omega 
\let\Y=\Upsilon


\definecolor{lightblue}{rgb}{0, 0.33, 0.71}



\DeclareMathOperator{\supp}{Supp}


\def\aa{\mathfrak{a}}
\def\bx{\bf x}
\def \blue#1 {\textcolor{blue}{#1}}
\def \red#1 {\textcolor{red}{#1}}
\def \blue#1{\textcolor{blue}{#1}}
%\overfullrule=5pt

%%%%%%%%%%%%%%%%%%%%%%%%%%%%%%%%%%%%
%workaround to use widecheck
\DeclareFontFamily{U}{mathx}{\hyphenchar\font45}
\DeclareFontShape{U}{mathx}{m}{n}{
      <5> <6> <7> <8> <9> <10>
      <10.95> <12> <14.4> <17.28> <20.74> <24.88>
      mathx10
      }{}
\DeclareSymbolFont{mathx}{U}{mathx}{m}{n}
\DeclareFontSubstitution{U}{mathx}{m}{n}
\DeclareMathAccent{\widecheck}{0}{mathx}{"71}

%%%%%%%%%%%%%%%%%%%%%%%%%%%%%%%%%%%%%%%%%%%%%%%%%%%%%%%%%%%%%%%%%%%%%%%%%%%%%%%%%%%%%%%%%%%%%%%%%%%%%%%%%%%%%%%%%%%%%%%%%%%%%%%%%%%%%%%%%%%%%%%%%%%%%%%%%%%%%%%%%%%%%%%%%%%%%%%%%%%%%%%%%%%%%%%%%%%%%%
\def\bskip{\\[-0.6cm]}


\title{Lee-Huang-Yang Energy for Dilute Bose Gases: \\ an Upper Bound for General Potentials} 

\date{\today}

\author{Giulia Basti\footnote{Department of Mathematics ``Guido Castelnuovo'', La Sapienza, Piazzale Aldo Moro, 5, 00185 Roma}, \, Morris Brooks\footnote{Institute of Mathematics, University of Zurich, Winterthurerstrasse 190, 8057 Zurich}, \, Serena Cenatiempo\footnote{Gran Sasso Science Institute, Viale Francesco Crispi 7, 67100 L'Aquila}, \, Alessandro Olgiati\footnote{Dipartimento di Matematica, Politecnico di Milano, Piazza Leonardo da Vinci 32, 20133 Milano}, \\ Benjamin Schlein$^\dagger$}


\begin{document}

\maketitle 

\begin{abstract}
We establish an upper bound for the ground state energy density of a dilute Bose gas, interacting through a general short-range, repulsive potential. To reach this goal, we extend the approach recently introduced in \cite{BBCOS} for hard spheres. 
\end{abstract}

\section{Introduction}

Based on the fundamental work of Bogoliubov \cite{B}, Lee-Huang-Yang \cite{LHY} considered a Bose gas at density $\rho > 0$, interacting through a two-body potential with scattering length $\frak{a} > 0$, and predicted that its ground state energy per unit volume has the form 
\begin{equation}\label{eq:LHY} 
e (\rho) = 4\pi \frak{a} \rho^2 \Big( 1+ \frac{128}{15 \sqrt{\pi}} (\rho \frak{a}^3)^{1/2} + \dots \Big) \end{equation} 
up to lower order terms, in the dilute limit $\rho \frak{a}^3 \ll 1$.

The first step towards a mathematical understanding of (\ref{eq:LHY}) is due by Dyson, who obtained a rigorous upper bound for $e(\rho)$, correct at leading order. The matching lower bound was derived by Lieb-Yngvason \cite{LY}. In the last years, the focus shifted to the second-order contribution on the r.h.s. of (\ref{eq:LHY}), known as the Lee-Huang-Yang term. In the so-called Gross-Pitaevskii regime, second order estimates for the ground state energy and for the low-energy excitation spectrum have been established in \cite{BBCS1,BBCS2} and, more recently, in \cite{HST,B}. In the thermodynamic limit, a lower bound capturing the Lee-Huang-Yang term has been shown for repulsive, quickly decaying and radial interactions, first in \cite{FS1}, under the additional assumption that the potential is integrable, and then in \cite{FS2}, in the general case. The first rigorous upper bound resolving the Lee-Huang-Yang term was proved in \cite{YY}, for sufficiently regular potentials. A simpler upper bound was later derived in \cite{BCS} for a larger class of repulsive, compactly supported and radial interactions; also this result, however, required $V \in L^3 (\bR^3)$ and excluded potentials with a hard core. For $V \in L^2 (\bR^3)$, a more precise upper bound, resolving also the third order term in the asymptotic expansion of $e(\rho)$, first predicted in the physics literature by Wu \cite{Wu}, Hugenholtz-Pines \cite{HP} and Sawada \cite{Sa}, was recently established in \cite{BOSS}. The corresponding lower bound is still open (but has been proven in the Gross-Pitaevskii regime in \cite{COSS}). New approaches have been developed in \cite{HHNST,HHST,FJGMOT} to extend (\ref{eq:LHY}) to an estimate for the free energy density, at low but positive temperatures.

Very recently, the restriction to integrable interactions appearing in \cite{YY,BCS} was lifted in \cite{BBCOS}, where an upper bound matching the Lee-Huang-Yang formula (\ref{eq:LHY}) was shown for interacting hard spheres. The goal of the present paper is to extend the analysis of \cite{BBCOS}, to derive an upper bound at Lee-Huang-Yang's order for general, repulsive, short-range interactions. 

We are going to consider a system of $N$ bosons moving in the box $\Lambda = [-L/2 ; L/2]^3$, with periodic boundary conditions and interacting through a repulsive potential $V : \mathbb{R}^3 \longrightarrow [0,\infty]$. We assume $V$ to be measurable, even and compactly supported, ie. $V(-x) = V(x)$ for all $x \in \bR^3$ and there exists $r_0 > 0$ with $\text{supp } V \subset B_{r_0} (0)$. Notice that our assumptions allow for potentials with a hard-core, ie. with $V(x) = \infty$, for all $|x| \leq r$, for some $r > 0$. The Hamilton operator has the form 
\begin{equation}\label{eq:HN0} H_N = \sum_{j=1}^N -\Delta_{x_j} + \sum^N_{i<j} V (x_i - x_j) \end{equation}
%and, in accordance with bosonic statistics, acts on %the Hilbert space $L^2_s (\Lambda^N)$, the subspace of %$L^2 (\Lambda^N)$ consisting of functions that are %symmetric with respect to permutations of the $N$ %particles. 
For a precise definition of $H_N$, we follow \cite{FS2} and introduce the quadratic form
\begin{align*}
    \mathcal{Q}_N [\Psi]:=\int_{\Lambda^N} \Big(\sum_{j=1}^N |\nabla_{x_j}\Psi|^2+\sum_{i<j}^N V(x_i-x_j)|\Psi|^2\Big) dx_1 \dots dx_N,
\end{align*}
acting on the domain 
\[ D(Q_N) = \Big\{ \Psi\in H_s^1(\Lambda^N) : \big(\sum_{i<j}^N V(x_i-x_j)\big)^{1/2}\Psi \in L^2(\Lambda^N) \Big\} \]
%, consisting of all  permutation symmetric $\Psi\in %H^1(\Lambda^N)$ such that $\big(\sum_{i<j} V(x_i-%x_j)\big)^{1/2}\Psi \in L^2(\Lambda^N)$.
where $H_s^1 (\Lambda^N)$ is the permutation symmetric subspace of $H^1 (\Lambda^N)$. On the Hilbert space 
\[ \cH_{\mathcal{Q}_N}:=\overline{ \{ \Psi\in L^2_s (\Lambda^N): \mathcal{Q}_N [\Psi]<\infty\}},\] the non-negative and closable form $\mathcal{Q}_N$ is densely defined; therefore the self-adjoint operator $H_N$ can be defined as its Friedrichs extension.

%At low density, the ground state energy depends on the scattering length of %the interaction. 
To introduce the scattering length of $V$, we consider the energy functional 
\begin{equation}\label{eq:Ef}
    \mathcal{E}[g]:=\int \left(2|\nabla g (x)|^2+V(x) |g (x)|^2 \right)dx .
\end{equation}
We denote by $f$ the unique minimizer of (\ref{eq:Ef}), on the space of all real-valued $g$, such that $g-1\in \dot{H}_1(\mathbb R^3)$; see \cite{RT}. The scattering length $\frak{a} \geq 0$ of $V$ is defined by 
\begin{equation}\label{eq:defa}
    8\pi \mathfrak{a}:= \mathcal{E} [f] = \inf_g \mathcal{E}[g].
\end{equation} 

We denote by $E_{N,L}$ the ground state energy of (\ref{eq:HN0}). We are interested in the ground state energy per unit volume, in the thermodynamic limit $N, L \to \infty$, at fixed density $\rho = N / L^3$, which is defined by 
\begin{equation}\label{eq:erho} 
e (\rho) = \lim_{N,L \to \infty : N/ L^3 = \rho} \frac{E_{N,L}}{L^3}  \, . 
\end{equation} 
The existence of the limit (\ref{eq:erho}) is well-known, see \cite{R}. Our main theorem provides an upper bound for (\ref{eq:erho}), in the dilute limit. 
\begin{theorem}\label{thm:main}
Let $V : \bR^3 \to [0;\infty]$ be measurable, even, with $\text{supp } V \subset B_{r_0} (0)$, for some $r_0 > 0$. Let $\frak{a} > 0$ be the scattering length of $V$, defined as in (\ref{eq:defa}). Moreover, let $e(\rho)$ denote the ground state energy density of the interacting Bose gas, defined in (\ref{eq:erho}). Then there exist $C > 0$ and a sufficiently small $\delta > 0$ such that 
\begin{equation}\label{eq:main} e(\rho) \leq 4\pi \frak{a} \rho^2 \Big( 1+ \frac{128}{15 \sqrt{\pi}} (\rho \frak{a}^3)^{1/2} + C (\rho \frak{a}^3)^{1/2+\delta} \Big) \end{equation} 
for all $\rho \frak{a}^3 > 0$ small enough. 
\end{theorem}

Our bound (\ref{eq:main}) is consistent with the Lee-Huang-Yang formula (\ref{eq:LHY}). The corresponding lower bound was established in \cite{FS2}, for measurable, spherical symmetric potentials $V : \bR^3 \to [0;\infty]$ decaying sufficiently fast at infinity. Combining Theorem \ref{thm:main} with the results of \cite{FS2}, we obtain a rigorous proof of the Lee-Huang-Yang formula, for all spherical symmetric, repulsive interactions with compact support. 


\medskip

{\it Acknowledgements.} G.B., S.C., and A.O. acknowledge funding from the European Research Council under the ERC Starting Grant MaTCh (Grant Agreement No. 101117299). G.B., S.C., and A.O. also thank GNFM–INdAM.
%G.B., S.C. and A.O. also warmly acknowledge the GNFM %(Gruppo Nazionale per la Fisica Matematica) - INDAM. 
A.O. acknowledges financial support from the MUR Grant “Dipartimento di Eccellenza 2023-2027” of Dipartimento di Matematica, Politecnico di Milano, and from the Politecnico di Milano ``Seed Fund Grant''. M.B. acknowledges funding from the UZH Postdoc Grant FK-25-103.


\section{The trial state}

To prove Theorem \ref{thm:main}, we are going to construct an appropriate trial state. To this end, it is convenient to work on the bosonic Fock space
\[ \cF = \bigoplus_{n=0}^\infty L^2 (\Lambda)^{\otimes_s n} \, . \]
On $\cF$, we introduce the usual creation and annihilation operators $a^* (f), a(f)$, defined for an arbitrary $f \in L^2 (\Lambda)$ by 
\[ \begin{split} (a (f) \Psi)^{(n)}  (x_1, \dots , x_n) &= \sqrt{n+1} \int_{\Lambda}  dx \bar{f} (x) \Psi^{(n+1)} (x, x_1, \dots , x_n) \\ 
(a^* (f) \Psi)^{(n)} (x_1, \dots, x_n) &= \frac{1}{\sqrt{n}} \sum_{j=1}^n f (x_j) \Psi^{(n-1)} (x_1, \dots , x_{j-1}, x_{j+1}, \dots , x_n)  \end{split} \]
and satisfying the canonical commutation relations 
\begin{equation*}\label{eq:CCR} 
\big[ a (f) , a^* (g) \big] = \langle f , g \rangle , \qquad \big[ a (f) , a(g) \big] = \big[ a^* (f) , a^* (g) \big] = 0  \,.
\end{equation*} 
For $p \in \Lambda^* = 2\pi \bZ^3/L$, we define the momentum-space creation and annihilation operators 
\[ \hat{a}_p^* = a^* (f_p), \quad \hat{a}_p = a (f_p) \]
with the normalized plane wave $f_p (x) = e^{ip\cdot x} / L^{3/2}$. In position space, we will use  operator-valued distributions $a_x^* , a_x$, defined for $x \in \Lambda$, so that 
\[ a^* (f) = \int f (x) a^*_x \, dx , \qquad a (f) = \int \bar{f} (x) \, a_x dx \]
for every $f \in L^2 (\Lambda)$. 

On $\cF$, we define the Hamilton operator $\cH = \cK + \cV$, where the kinetic energy operator $\cK$ can be written in momentum space as 
\begin{equation}\label{eq:cK} \cK = \sum_{p \in \Lambda^*} p^2 \hat{a}_p^* \hat{a}_p \end{equation}
while the potential energy operator $\cV$ is conveniently expressed in position space as 
\begin{equation}\label{eq:cV} \cV = \frac{1}{2} \int dx dy \, V(x-y) a_x^* a_y^* a_y a_x  \,. \end{equation} 
Our goal is to construct a normalized trial state $\Psi \in \cF$, with approximately $N = \rho L^3$ particles and with energy per unit volume $\langle \Psi, \cH \Psi \rangle / L^3$ sufficiently close to the r.h.s. of (\ref{eq:main}), which will then follow by a standard equivalence of ensembles argument. To define such a trial state, we have to generate, first of all, a Bose-Einstein condensate in the zero-momentum state, with density $0 < \rho_0 \leq \rho$. To reach this goal, we will use the unitary Weyl operator 
\begin{equation}\label{eq:weyl} W(\rho_0) = e^{a^* (\sqrt{\rho_0}) - a (\sqrt{\rho_0})} = e^{\sqrt{\rho_0} \int dx (a_x^* - a_x)} \end{equation} 
producing the coherent state 
\begin{equation} \label{eq:coh} W(\rho_0) \Omega = e^{-\rho_0/2} \Big\{ 1, \sqrt{\rho_0}, \dots, \frac{\rho^{k/2}_0}{\sqrt{k!}}, \dots \Big\} \in \cF \,.\end{equation}
Energetically, choosing $\rho_0 = \rho$, the pure condensate (\ref{eq:coh}) is very far from the true ground state energy, because it completely lacks correlations. In fact, if the potential $V$ has a hard-core, the energy of (\ref{eq:coh}) is even infinite. Hence, we have to modify (\ref{eq:weyl}), to introduce the correct correlation structure, making sure, in particular, that the trial state vanishes on hard-cores. To this end we proceed as in \cite{BBCOS}. We introduce two length scales $\ell = \frak{a} (\rho \frak{a}^3)^{-\delta}$ and $\ell_0 = (\rho \frak{a})^{-1/2} (\rho \frak{a}^3)^{-\eps}$, for $\delta , \eps > 0$ small enough, to be fixed later on. In the dilute limit $\rho \frak{a}^3 \ll 1$, $\ell$ is a bit larger than the range of the interaction while $\ell_0$ is just a bit larger than the healing length $(\rho \frak{a})^{-1/2}$ of the gas, ie. $\frak{a} \leq r_0 \ll \ell \ll (\rho \frak{a})^{-1/2} \ll \ell_0$. We will modify the coherent state (\ref{eq:coh}) adding a Jastrow factor acting on short length scales $|x| \leq \ell$ and a Bogoliubov transformation generating correlations on larger scales, $\ell \leq |x| \leq \ell_0$. The Jastrow factor is crucial for the trial state to vanish on hard-cores. The Bogoliubov transformation, on the other hand, is crucial to describe correlations up to (and just beyond) the healing length, as needed to reach Lee-Huang-Yang precision. 

To define the Jastrow factor on short length scales, we need some properties of the minimizer of the energy functional (\ref{eq:Ef}), defining the scattering length of the interaction potential $V$.
\begin{lemma}
\label{Lem:general_potential_eff}
Assume that $V : \bR^3 \to [0;\infty]$ is measurable, even, with $\text{supp } V \subset B_{r_0} (0)$, for some $r_0 > 0$. Let $f$ denote the unique minimizer of (\ref{eq:Ef}), on the space of all $f$ such that $f-1\in \dot{H}_1(\mathbb R^3)$. We consider the distribution 
\begin{equation}\label{eq:Veff} V_\mathrm{eff} = 2\Delta f \, . \end{equation}
Then, $V_\mathrm{eff}$ is non-negative, with support in $B_{r_0}(0)$. Its Fourier transform $\widehat{V}_\mathrm{eff} \in L^\infty (\bR^3)$ is such that $\widehat{V}_\mathrm{eff} (0) = 8\pi \frak{a}$ and, for every $m \in \bN$, $\| \nabla^m \widehat{V}_\mathrm{eff} \|_\infty < \infty$. Furthermore, for $|x|> r_0$, we have 
    \begin{align}
    \label{Eq:formula_f_integ}
        f(x)=1-\int \frac{V_\mathrm{eff}(y)}{8\pi |x-y|} dy\,.
    \end{align}
    In particular, it follows that 
    \begin{equation} \label{eq:1-f} 
    0 \leq 1- f (x) \lesssim \frac{\frak{a}}{|x|} \end{equation} 
    for all $x \in \bR^3$ and 
    \begin{equation}\label{eq:nablaf} |\nabla f (x)| \lesssim \frac{\frak{a}}{x^2} \end{equation} 
for all $|x| > 2 r_0$. \end{lemma}

{\it Remark.} Due to the compact support of $V_\mathrm{eff}$, both the Fourier transform
\begin{align*}
   \widehat{V}_\mathrm{eff}(k) = \int e^{-iky}V_\mathrm{eff}(y) dy := \int e^{-iky} \varphi(y)V_\mathrm{eff}(y) dx
\end{align*}
as well as the expression 
\begin{align}
\label{Eq:extension_to_bounded_func}
    \int \frac{V_\mathrm{eff}(y)}{8\pi |x-y|} dy :=  \int \frac{\varphi(y)V_\mathrm{eff}(y)}{8\pi |x-y|} dy
\end{align}
are well-defined, for $|x| > r_0$, if we choose $\varphi$ a smooth and compactly supported test function with $\varphi=1$ on $B_{r'}(0)$ and $\varphi=0$ on $\mathbb R^3\setminus B_{r''}(0)$ for some $r_0 < r' < r'' < |x|$. Notably, these definitions are independent of the concrete choice of $\varphi$.

{\it Remark.} If we additionally assumed that $V$ is radial, we would find $f(x) = 1 - \frak{a} / |x|$, for $|x| > r_0$. This would immediately imply the bounds (\ref{eq:1-f}), (\ref{eq:nablaf}). 

\begin{proof}[Proof of Lemma \ref{Lem:general_potential_eff}]
From the variational definition, using test functions of the form $f\pm \epsilon g$ and sending $\epsilon\rightarrow 0$, it is immediately clear that $V_\mathrm{eff}$ is supported on $B_{r_0}(0)$, i.e.
\begin{align}
\label{Eq:Scattering_I}
    \langle V_\mathrm{eff}, g\rangle=0,
\end{align}
if $g(x)=0$ for all $|x|\leq r_0$. Similarly, using test functions of the form $f+ \epsilon g$, $\epsilon > 0$, we conclude that $V_\mathrm{eff}$ is non-negative, i.e.
\begin{align*}
     \langle V_\mathrm{eff}, g\rangle\geq 0
\end{align*}
in case $g\geq 0$. To prove (\ref{Eq:formula_f_integ}), let $m_\epsilon(y):=\epsilon^{-3}m(\epsilon^{-1} y)$, for $\epsilon > 0$ and for a smooth function $m$, supported in $B_1 (0)$ and with $\int m(y) \mathrm{d}y=1$. For every $\epsilon > 0$, $m_\epsilon * V_\text{eff}$ is a smooth function, supported in $B_{r_0+\epsilon} (0)$. Choosing $\varphi$ as in \eqref{Eq:extension_to_bounded_func} and $\epsilon > 0$ so small that $r_0 + \epsilon < r'$, we conclude that 
%for a given $|x|>r_0$. Since $y\mapsto \frac{\varphi(y)}{8\pi |x-y|}$ is an %admissible test function, we have
\begin{align*}
    \int \frac{V_\mathrm{eff}(y)}{8\pi |x-y|} dy=\lim_{\epsilon\rightarrow 0}\int \frac{(m_\epsilon* V_\mathrm{eff})(y)}{8\pi |x-y|} dy \,.
\end{align*}
Consequently, to verify \eqref{Eq:formula_f_integ}, it is enough to show that, for all $\epsilon>0$ and $x \in \bR^3$,
\begin{align}
\label{Eq:formula_f_integ_in_proof}
    m_\epsilon*(f-1) (x) =-\int \frac{(m_\epsilon* V_\mathrm{eff})(y)}{8\pi |x-y|} dy \,.
\end{align}
By the definition of $V_\mathrm{eff}$, it is clear that $2\Delta m_\epsilon*(f-1) = m_\epsilon * V_\mathrm{eff}$. Furthermore, the smooth function on the right hand side of \eqref{Eq:formula_f_integ_in_proof} satisfies 
\begin{align*}
    2\Delta \left(-\int \frac{(m_\epsilon* V_\mathrm{eff})(y)}{8\pi |x-y|} dy\right)= (m_\epsilon * V_\mathrm{eff}) (x)
\end{align*}
as well. This concludes the proof of \eqref{Eq:formula_f_integ_in_proof}, as both the function on the left hand side as well as the one on the right hand side of \eqref{Eq:formula_f_integ_in_proof} vanish at infinity.  
%(by the maximum principle, an harmonic function vanishing as $|x| \to \infty$ %vanishes everywhere). 

In order to verify that $\widehat{V}_\text{eff} (0) = 8\pi \frak{a}$ we observe, first of all, that 
\begin{equation}\label{eq:Veff0} \widehat{V}_\text{eff} (0) = \lim_{R \to \infty} 2 \int_{B_R (0)} \Delta f \, dx = \lim_{R \to \infty} 2 \int_{\partial B_R (0)} \nabla f  \cdot n \, d\sigma \, . \end{equation} 
On the other hand, using the variational definition of $\frak{a}$ and observing that $f \pm \epsilon \varphi f$ is an admissible test function, for a smooth and compactly supported $\ph$, with $\ph = 1$ on the support of $V$, we obtain 
\begin{align*}
    0&=\int \left[2 \nabla f \cdot  \nabla (\varphi f)+V f (\varphi f)\right] dx \\ &=\int \left[2 \nabla f \cdot  \nabla (\varphi f)+V f^2\right] dx \, = 8\pi \frak{a} + \int 2 \, \nabla f \cdot \nabla (\ph f - f)  \, dx 
\end{align*}
implying 
\[ \int 2 \, \nabla f \cdot \nabla (f-\ph f) \, dx = 8\pi \frak{a} \, . \]
Since furthermore $-\Delta f=0$ on the support of $1-\varphi$, we conclude that 
\[ \begin{split} \int 2 \, \nabla f \cdot \nabla (f-\ph f) \, dx &= \lim_{R \to \infty} \int_{B_R (0)} 2 \, \nabla f \cdot \nabla (f-\ph f) \, dx \\ &= \lim_{R \to \infty}
\int_{\partial B_R (0)} 2\nabla f \cdot n  \, (f - \ph f)  \, d\sigma \\ &= \widehat{V}_\text{eff} (0) + \lim_{R \to \infty} \int_{\partial B_R (0)}  2\nabla f \cdot n  \, (f-1) \, d\sigma \end{split} \]
where in the last line we used (\ref{eq:Veff0}) and the fact that $\ph =0$ on $\partial B_R (0)$, for $R$ large enough. To show that $\widehat{V}_\text{eff} (0) = 8\pi \frak{a}$, we just observe that, from \eqref{Eq:formula_f_integ}, $|f (x)-1| \lesssim 1/R$, $|\nabla f (x)| \lesssim 1/R^2$ for $x \in \partial B_R (0)$ and $R$ large enough, implying that the integral on the r.h.s. vanishes, as $R \to \infty$. From $\widehat{V}_\text{eff} (0) = 8\pi \frak{a}$ we also obtain, since $V_\mathrm{eff} \geq 0$, that $\| V_\mathrm{eff} \|_1 = \sup \langle V_\mathrm{eff} , g \rangle = 8\pi \frak{a}$, where the supremum is taken over all test functions $g$, with $\| g \|_\infty \leq 1$. Similarly, as a consequence of \eqref{Eq:Scattering_I}, the $L^1$-norm of $x_{i_1}\dots x_{i_m}V_\mathrm{eff}(x)$ is finite as well, and therefore 
\begin{align*}
    \|\partial_{k_{i_1}}\dots \partial_{k_{i_m}} \widehat{V}_\mathrm{eff}\|_\infty < \infty \,.
\end{align*}

Finally, to show (\ref{eq:1-f}), (\ref{eq:nablaf}), we use \eqref{Eq:formula_f_integ} and we remark that, for $|x| > 2 r_0$ and $|y| \leq r_0$, $|x-y| \geq |x|/2$. Together with the identity $\widehat{V}_\text{eff} (0) = 8\pi \frak{a}$, we immediately find the desired bounds, for $|x| \geq 2 r_0$. Since $0 \leq 1- f(x) \leq 1$ everywhere (by the variational characterization of $f$), (\ref{eq:1-f}) can be extended to all $x \in \bR^3$.  
\end{proof}

Next, we truncate the minimizer $f$ of the energy functional (\ref{eq:Ef}) at the length scale $\ell = \frak{a} (\rho \frak{a}^3)^{-\delta}$, defining 
\begin{equation}\label{eq:fell} f_\ell (x) = \chi_\ell (x) f(x) + (1- \chi_\ell (x)) = 1- \chi_\ell (x) (1 -f(x)) \end{equation}
where $\chi_\ell (x) = \chi (x/\ell)$ and $\chi \in C^\infty_c (\bR^3)$, with $\chi (x) = 1$ for $|x| \leq 2$, $\chi (x) = 0$, for $|x| > 4$. 
\begin{lemma}  \label{lm:omega}
We have $0\leq f_\ell (x) \leq 1$ for all $x \in \bR^3$. Let $\omega_\ell (x) = 1- f_\ell (x)$. Then 
\[ \omega_\ell (x) \lesssim \frac{\frak{a}}{|x|} \chi_\ell (x) \]
for all $x \in \bR^3$ and 
\[ |\nabla f_\ell (x)| = |\nabla \omega_\ell (x)| \lesssim \frac{\frak{a}}{x^2} \]
for all $|x| > 2 r_0$. Moreover, 
\begin{equation}\label{eq:en-fl} \Big| \int \big[ |\nabla f_\ell (x)|^2 + \frac{1}{2} V (x) |f_\ell (x)|^2 \big] dx - 4 \pi \frak{a} \Big| \lesssim \frac{\frak{a}^2}{\ell} \,.\end{equation}
\end{lemma} 
\begin{proof}
The bounds for $f_\ell, \omega_\ell, \nabla f_\ell$ follow from Lemma \ref{Lem:general_potential_eff}. To show (\ref{eq:en-fl}), we recall the definition (\ref{eq:defa}) and we observe that, since $f_\ell (x) = f(x)$ on the support of $V$ and since $\nabla f_\ell (x) = \chi_\ell (x)\nabla f (x)  - (1-f(x)) \nabla \chi_\ell (x)$, 
\[ \begin{split} \Big| \int  \big[ |\nabla f_\ell (x)|^2 &+ \frac{1}{2} V (x) | f_\ell (x)|^2\big] dx  - 4\pi \frak{a} \Big| \\ \leq \; &\int |\nabla f (x)|^2 (1- \chi_\ell^2 (x)) dx + \int |\nabla \chi_\ell (x)|^2 (1- f(x))^2 \\ &+2 \int |\nabla f (x)| |\nabla \chi_\ell (x)| (1- f(x)) dx \,. \end{split} \]
From Lemma \ref{Lem:general_potential_eff}, we have $|1-f(x)| \lesssim \frak{a}/|x|$ and $|\nabla f (x)| \lesssim \frak{a}/x^2$ on the supports of $1-\chi_\ell^2$ and of $\nabla \chi_\ell$. With $|\nabla \chi_\ell (x)| \lesssim \mathbbm{1} (2\ell \leq |x| \leq 4\ell) / \ell$, we obtain (\ref{eq:en-fl}). 
\end{proof}

With $f_\ell$ as in (\ref{eq:fell}), we define $J : \cF \to \cF$, requiring that  
\begin{equation}\label{eq:defJ} (J \Phi)^{(n)} (x_1, \dots , x_n) = \prod_{i<j}^n f_\ell (x_i - x_j) \Phi^{(n)} (x_1, \dots , x_n) \,. \end{equation}
For $x \in \Lambda$, we also define $J(x) : \cF \to \cF$ by 
\[ (J(x) \Phi)^{(n)} (x_1, \dots , x_n) = \prod_{j=1}^n f_\ell (x - x_j) \Phi^{(n)} (x_1, \dots , x_n) \,.\]
Then, we have the identities 
\[ a_x J = J(x) J a_x , \qquad J a_x^* = a_x^* J(x) J \]
and
\[ a_y J(x) = f_\ell (x-y) J(x) a_y, \qquad J(x) a_y^* = f_\ell (x-y) a_y^* J(x) \,.\]
Clearly, $0 \leq J(x), J \leq 1$ for all $x \in \bR^3$. From
\[ 1 - \prod_{j=1}^n f_\ell (x- x_j) \leq \sum_{j=1}^n (1 - f_\ell (x-x_j)) = \sum_{j=1}^n \omega_\ell (x-x_j) \]
we also find 
\begin{equation}\label{eq:1-J} 1 - J(x) \leq d\Gamma (\omega_{\ell,x}) = \int dy \, \omega_\ell (x-y) a_y^* a_y \,. \end{equation}

Acting with (\ref{eq:defJ}) on the condensate (\ref{eq:coh}) we generate a Jastrow factor, capturing correlations up to the length scale $\ell$ and making sure that the trial state vanishes on hard-cores (if $V (x) = \infty$ for all $|x| \leq r$, then clearly the minimizer $f$ of (\ref{eq:Ef}) satisfies $f(x) = 0$ for all $|x| \leq r$). To create correlations on larger length scales, we are going to use a Bogoliubov transformation. Proceeding as in \cite{BBCOS}, we define the coefficients \begin{equation}\label{eq:hatsk} \begin{split} 
\widehat{s}_p &= - \frac{p^2 + \rho_0 \widehat{V}_\text{eff} (p) -  \sqrt{|p|^4 + 2p^2 \rho_0 \widehat{V}_\text{eff} (p)}}{\sqrt{\rho_0^2 \widehat{V}_\text{eff} (p)^2 - \Big( p^2 + \rho_0 \widehat{V}_\text{eff} (p) -  \sqrt{|p|^4 + 2p^2 \rho_0 \widehat{V}_\text{eff} (p)}\Big)^2}}
\end{split} 
\end{equation}
for every $p \in \Lambda^* \backslash \{ 0 \}$, with $V_\mathrm{eff}$ as in (\ref{eq:Veff}). We set also $\widehat{s} (0) = 0$ and we denote by $s(x)$ the periodic function on $\Lambda$, with Fourier coefficients (\ref{eq:hatsk}). Then, we define  
\begin{equation}\label{eq:tlsigma} \tilde{\sigma} (x) = \frac{\chi_{\ell_0} (2x) (1- \chi_\ell (2x))}{f_\ell (x)} \big[ \rho_0 \omega_\ell (x) + s (x) \big] \end{equation} 
and 
\begin{equation}
\label{eq:def-sigma} \sigma (x) = \tilde{\sigma} (x) - \ell_0^{-3} \Big[ \int \tilde{\sigma} (z) dz \Big]  \ph (x/\ell_0) 
\end{equation} 
so that $\widehat{\sigma}_0 = 0$. We set now $\widehat{\eta}_p = \mathrm{sinh}^{-1} (\widehat{\sigma}_p)$, with $\widehat{\sigma}_p$ denoting the Fourier coefficients of (\ref{eq:def-sigma}), and we consider the Bogoliubov transformation 
\begin{equation}\label{eq:def-bog} T = \exp \Big[ \frac{1}{2} \sum_{p \in \Lambda^* \backslash \{ 0 \}} \widehat{\eta}_p \big( \hat{a}_p^* \hat{a}_{-p}^* - \hat{a}_p \hat{a}_{-p} \big) \Big] \,. \end{equation} 
The action of $T$ on creation and annihilation operators is explicitly given by 
\[ T^* \hat{a}_p T = \widehat{\gamma}_p \hat{a}_p + \widehat{\sigma}_p \hat{a}_{-p}^* \]
for every $p \in \Lambda^*$, with $\widehat{\gamma}_p = \cosh \eta_p = \cosh (\sinh^{-1} (\widehat{\sigma}_p))$. In the next lemma, which is an adaptation of \cite[Lemma 2.2]{BBCOS}, we collect some important properties of the kernels $\eta, \sigma, \gamma$. We briefly discuss its proof (focusing, in particular, on the differences w.r.t. the proof of \cite[Lemma 2.2]{BBCOS}) in Appendix \ref{app:kernels}. 
\begin{lemma} \label{lm:eta}
The coefficients $\widehat{s}_k$, defined in \eqref{eq:hatsk}, satisfy 
\begin{equation} \label{eq:sfourier} |\widehat{s}_k| \leq C \min \Big\{  \frac{\rho^{1/4}}{|k|^{1/2}} , \frac{\rho}{k^2} \Big\} \end{equation} 
for all $k \in \Lambda^*_+ = \Lambda^* \backslash \{ 0 \}$. For the function $s : \Lambda \to \bR$ with Fourier coefficients $\widehat{s}_k$ we find the pointwise bounds 
\begin{equation*} | s (x)| \leq C \frac{\rho}{|x|} \min \Big\{ 1 , \frac{1}{(\rho^{1/2} |x|)^{3/2}} \Big\} , \qquad |\nabla s (x)| \leq C \frac{\rho |\hspace{-.05cm} \log \rho |}{x^2} \end{equation*} 
for all $|x| \geq 15r_0$. Moreover 
\begin{equation}\label{eq:est-combi}
 \int_{|x|\geq \ell_0} |s(x)|^2 dx  \leq C\ell_0^{-2}\rho^{\frac{1}{2}}  , \;   \int_{2\frak{a} \leq |x|\leq  2\ell} 
 |s(x)|^2 dx  \leq C\ell\rho^{2}, \;  \int_{|x|\geq \ell_0} |\nabla s(x)|^2 dx  \leq  C\ell_0^{-4}\rho^{\frac{1}{2}}\, . \end{equation} 
 %
 %
%\begin{align}
 %\label{Th:est_1}
%        \int_{|x|\geq \ell_0} |s(x)|^2\mathrm{d}x & \leq C\ell_0^{-2}\rho^{\frac{1}{2}},\\
%         \label{Th:est_1_1}
%        \int_{|x|\leq  2\ell} |s(x)|^2\mathrm{d}x & \leq C\ell\rho^{2},\\
%            \label{Th:est_2}
%        \int_{|x|\geq \ell_0} |\nabla s(x)|^2\mathrm{d}x &  \leq  C\ell_0^{-4}\rho^{\frac{1}{2}}\,.
%\end{align}
%
Let now $\zeta\in \{\tilde\sigma, \sigma,\gamma-\mathbbm{1},\mathbbm{1} - \gamma^{-1} , \gamma^{-1}*\sigma\}$. Then we have 
\begin{equation} \label{eq:zeta-combi1} 
 \|\zeta\|^2_2 \leq C \rho^{\frac{3}{2}}, \qquad \| \zeta \|_\infty \leq C\, \frac{\rho}{\ell} , \qquad  \|\zeta\|_1  \leq C \big(\rho^{\frac{1}{2}}\ell_0\big)^3 
 \end{equation} 
 and also 
 \begin{equation}\label{eq:zeta-combi2}
   \|\nabla \zeta\|^2_2 \leq C \rho^{2} , \qquad  \|\nabla \zeta\|_\io  \leq C \, \frac{\rho}{\ell^2}
   \end{equation} 
%
%  \begin{align}
%            \label{Th:L_2_est}
%        \|\zeta\|^2_2& \leq C \rho^{\frac{3}{2}},\\
%          \label{Th:est_3_alt}
%             \|\zeta\|_1 & \leq C \big(\rho^{\frac{1}{2}}\ell_0\big)^3,\\
%\label{Th:est_3b} 
%\| \zeta \|_\infty &\leq C \rho / \ell \\ 
%  \label{Th:H_1_est}
%        \|\nabla \zeta\|^2_2& \leq C \rho^{2}  \\
%          \label{Th:est_4_alt_2}
% \|\nabla \zeta\|_\io & \leq C \frac{ \rho}{\ell^2}\ .
%\end{align} 
and, for any $m \in \bN$, the pointwise decay estimate 
\begin{equation}
            \label{eq:decay}
 |\zeta(x)|  \leq C \, \frac{\rho}{\ell}\Big( \frac{\rho^{\frac{1}{4}}\ell_0^{\frac{3}{2}}}{|x|}\Big)^m \leq C \, \frac{\rho}{\ell} \, \frac{1}{\big( \rho^{1/2+3\eps/2} |x| \big)^{m}} \,.
\end{equation}
For $\zeta \in \{ \tilde{\sigma}, \sigma \}$, we get stronger estimates. In particular
\begin{equation}\label{eq:s2-restricted}
\int_{|x|\geq \ell_0} \big[ |\tilde{\sigma} (x)|^2 + |\sigma (x)|^2 \big] dx  \leq C\ell_0^{-2}\rho^{\frac{1}{2}}  , \qquad  \int_{|x|\leq  2\ell} 
 \big[ |\tilde{\sigma} (x)|^2 + | \sigma (x)|^2 \big] dx \leq C\ell\rho^{2}
 \end{equation} 
and  
\begin{equation} 
            \label{Th:est_3}
             \|\sigma\|_1  \leq 2\|\tilde \sigma\|_1\leq C \big( \rho^{\frac{1}{2}}\ell_0\big)^\frac{1}{2}\,.
             \end{equation}
Additionally, we get the pointwise bounds 
\begin{equation}\label{eq:snabla-point} |\nabla \tilde{\sigma} (x)| , |\nabla \sigma (x) | \leq C \frac{\rho | \hspace{-.03cm} \log \rho|}{x^2} \end{equation} 
for all $x \in \Lambda$, which imply  
\begin{equation}\label{eq:snabla-p}\| \nabla \sigma \|_p, \| \nabla \tilde{\sigma} \|_p  \leq C \rho |\hspace{-.03cm} \log \rho  | \left\{ \begin{array}{ll}  \ell^{3/p-2} \quad \text{if } p > 3/2 \\ \ell_0^{3/p-2} \quad \text{if } p < 3/2 \end{array} \right.\end{equation} 
and
\begin{equation}\label{eq:nu-nabla-p} \| \nabla (\gamma^{-1}*\sigma) \|_p  \leq C (\rho^{1/2} \ell_0 )^3 \rho |\hspace{-.03cm} \log \rho | \left\{ \begin{array}{ll}  \ell^{3/p-2} \quad &\text{if } p > 3/2 \\  \ell_0^{3/p-2} \quad &\text{if } p < 3/2 \end{array} \right.\end{equation} 
for all $1 \leq p \leq \infty$.
\end{lemma}

We are going to consider the trial state 
\begin{equation}\label{eq:trial} \Psi = \frac{1}{Z} J W(\rho_0) T \Omega \in \cF \end{equation}
with $J$ as in (\ref{eq:defJ}), the Weyl operator $W(\rho_0)$ as in (\ref{eq:weyl}), the Bogoliubov transformation $T$ from (\ref{eq:def-bog}) and with the normalization constant $Z = \| J W(\rho_0) T \Omega \|$. We choose $0 < \rho_0 < \rho$ such that 
\begin{equation}\label{eq:choice} \rho+ C _1 \rho^{7/4-11\eps-\delta} \leq \rho_0 + \| \sigma \|^2 \leq \rho + C_2 \rho^{7/4-11\eps-\delta} \end{equation} 
for fixed constants $C_1 < C_2$. This is possible because  $\| \sigma \|_2^2 \simeq \rho^{3/2} \gg \rho^{7/4-11\eps-\delta}$.

As observed in \cite[Lemma 2.3]{BBCOS}, we have the important identity 
\begin{equation}\label{eq:id} a_x \Psi = \sqrt{\rho_0}J(x)\Psi+J(x)\int dy\,(\gamma^{-1} * \sigma)(x-y)a^*_y J(y)\Psi.
\end{equation}

The proof of Theorem \ref{thm:main} is based on the following two propositions, estimating number of particles and energy of the trial state (\ref{eq:trial}). 
\begin{proposition} \label{prop:N_on_psi}
	We have 
	\begin{equation} \label{eq:cN_ub}
		\big| \langle \Psi,\mathcal{N}\Psi\rangle - \rho L^3 \big| \leq C \rho^{7/4-11\eps-\delta} L^3
	\end{equation}    
	if $\eps , \delta > 0$ are small enough.  	
\end{proposition}

\begin{proposition} \label{prop:Psi-energy} 
%Let $\Psi$ be defined as in \eqref{eq:trial} with $\ell= \aa (\rho\aa^3)^{-\d}$ and $\ell_0=(\rho \aa)^{-1/2}(\rho\aa^3)^{-\eps}$. 
Let $0 < \delta < \eps/2$, with $\eps > 0$ small enough. Then there exist $C>0$ such that
	\begin{equation} \label{eq:Hpsi0}
		L^{-3} \langle \Psi,\mathcal{H} \Psi\rangle \leq 4 \pi \frak{a} \rho^2 \Big(1 + \frac{128}{15\sqrt{\pi}} (\rho \frak{a}^3)^{1/2} + C (\rho \frak{a}^3)^{1/2+\delta} \Big) 
	\end{equation}
	for all $\rho \frak{a}^3 > 0$ small enough. 
\end{proposition}

\begin{proof}[Proof of Theorem \ref{thm:main}]
Given Prop. \ref{prop:N_on_psi} and Prop. \ref{prop:Psi-energy}, the proof of Theorem \ref{thm:main} is based on equivalence of ensembles. The argument is identical to the one used in \cite{BBCOS}. 
\end{proof}


\section{Number of particles in the trial state}

In this section, we show Prop. \ref{prop:N_on_psi}. To this end, we recall from \cite{BBCOS} that (\ref{eq:trial}) satisfies some important a-priori bounds on the local number of particles and also on the local number of excitations of the Bose-Einstein condensate, as measured by the fields 
\begin{equation}\label{eq:bb} \begin{split} b^*_x &= a^*_x - \sqrt{\rho_0} = W (\rho_0) a^*_x W(\rho_0)^* \\  b_x &= a_x - \sqrt{\rho_0} = W(\rho_0) a_x W(\rho_0)^*\,. \end{split} \end{equation} 
The following lemma is shown in \cite[Lemma 4.3, Lemma 4.4]{BBCOS}.
\begin{lemma} \label{lm:a-bds} 
Let $R = C \rho^{-1/2-3\eps}$, $k \in \bN \backslash \{ 0 \}$. Then we have 
\begin{equation}\label{eq:kas}  \int dx_1 \dots dx_k \, \prod_{j=2}^k \chi (|x_1 - x_j| \leq R)  \langle \Psi, a_{x_1}^* \dots a_{x_k}^* a_{x_k} \dots a_{x_1} \Psi \rangle \leq C \rho^{-(k-3)/2 - (12 k -9) \eps} L^3 \end{equation} 
and 
\begin{equation}\label{eq:kbs}  \int dx_1 \dots dx_k \, \prod_{j=2}^k \chi (|x_1 - x_j| \leq R)  \langle \Psi, b_{x_1}^* \dots b_{x_k}^* b_{x_k} \dots b_{x_1} \Psi \rangle \leq C \rho^{3/2 - (10k -9) \eps} L^3 \end{equation} 
if $\eps, \delta > 0$ and $\rho > 0$ are small enough (if $k=1$, we remove $\prod_{j=2}^k \chi (|x_1 - x_j| \leq R)$). Moreover, if we have only one annihilation (or creation) operator, we find \begin{equation}\label{eq:1bb} \Big| \sqrt{\rho_0} \int dx\langle \Psi, a_x \Psi \rangle - \rho_0 L^3 \Big| = \Big| \sqrt{\rho_0} \int dx \langle \Psi, b_x \Psi \rangle \Big| \lesssim \rho^{2-9\eps-2\delta} L^3 \end{equation} 
if $\eps, \delta > 0$ and $\rho >0$ are small enough. 
\end{lemma} 

Sometimes, it is useful to estimate the number of excitations generated only by the Jastrow factor (removing excitations created by the Bogoliubov transformation). Defining the new operators
\begin{equation}\label{eq:def-cc}
\begin{split} 
c^*_x &=b^*(\g_x)-b(\s_x)= a^* (\gamma_x) - a (\sigma_x) - \sqrt{\rho_0} = T W (\rho_0) a^*_x W^* (\rho_0) T^* \\
c_x &=b(\g_x)-b^*(\s_x)= a (\gamma_x) - a^* (\sigma_x) - \sqrt{\rho_0} =  T W (\rho_0) a_x W^* (\rho_0) T^* 
\end{split} \end{equation} 
we have the following bounds, established in \cite[Lemma 4.5]{BBCOS}.
\begin{lemma} \label{lm:c's}
We have 
\begin{equation}\label{eq:cc-claim} \int dx \langle \Psi, c_x^* c_x \Psi \rangle \leq C \rho^{2-16\eps-2\delta} L^3\,. \end{equation} 
Moreover, let $R = C \rho^{-1/2-3\eps}$. Then, we have 
\begin{equation}\label{eq:cccc-claim} \int_{|x-y| \leq R} dx dy  \, \langle \Psi,  c_x^* c_y^* c_y c_x  \Psi \rangle \leq C \rho^{2-30\eps-\delta} L^3 \,.
\end{equation} 
More generally, for $\ph, \tau \in \{ \sigma, \gamma^{-1}\ast \sigma, \gamma -\mathbbm{1}, \mathbbm{1} \}$, we find (if $\tau = \mathbbm{1}$, we set $c (\tau_x) = c_x$) 
\begin{equation}\label{eq:cc-gh}
\int dx \, \langle \Psi, c^* (\tau_x) c (\tau_x) \Psi \rangle \leq C \rho^{2-22\eps-2\delta} L^3 \end{equation}
and  
\begin{equation}\label{eq:cccc-gh} 
\int_{|x-y| \leq R} dx dy  \, \langle \Psi,  c^* (\tau_x) c^* (\ph_y) c (\ph_y) c (\tau_x)  \Psi \rangle \leq C \rho^{2-42\eps-\delta} L^3 \,.
\end{equation} 
\end{lemma} 
As an application of the last lemma, we can show stronger estimates for the local number of particles and of excitations, in small regions. The next lemma is taken from \cite[Lemma 4.6]{BBCOS} .
\begin{lemma} \label{lm:aaaaC} 
For $0 < \delta < \eps$, $\eps > 0$ small enough and $\rho > 0$ small enough, we have 
\begin{equation}\label{eq:aaaaC} \int_{|x-y| < C \ell} dx dy \, \langle \psi, b_x^* b_y^* b_y b_x \psi \rangle , \; \int_{|x-y| < C \ell} dx dy \, \langle \psi, a_x^* a_y^* a_y a_x \psi \rangle  \lesssim \rho^{2-36\eps-\delta} L^3\,. \end{equation} 
Moreover, recalling $R = C \rho^{-1/2-3\eps}$, we find  
\begin{equation}\label{eq:aaaCR} \int_{|x-y| \leq C \ell, |x-z| \leq R} dx dy dz \langle \Psi, a_x^* a_y^* a_z^* a_z a_y a_z \Psi \rangle \lesssim \rho^{3/2-45\eps-\delta} L^3\,. \end{equation}
\end{lemma} 

We can now prove Prop. \ref{prop:N_on_psi}, estimating the total number of particles in (\ref{eq:trial}). Also this proof is taken from \cite{BBCOS}. 
\begin{proof}[Proof of Prop. \ref{prop:N_on_psi}] 
Writing  
\begin{equation} \label{eq:a_to_b}
a_x = \sqrt{\rho_0} + b_x\,, \qquad a_x^* = \sqrt{\rho_0} + b_x^*
\end{equation} 
and 
\[ b_x = c (\gamma_x) + c^* (\sigma_x), \qquad b_x^* = c^* (\gamma_x) + c (\sigma_x) \]
we obtain 
\[ \begin{split}  \langle\Psi ,\mathcal{N} \Psi\rangle =\; &\rho_0 L^3 +  2\sqrt{\rho_0} \, \text{Re } \int dx \langle \Psi, b_x \Psi \rangle + \int dx \langle \Psi, b_x^* b_x \Psi \rangle 
\\ =\;&(\rho_0 + \| \sigma \|_2^2) L^3 +  2 \sqrt{\rho_0} \, \text{Re } \int dx \langle \Psi, b_x \Psi \rangle \\ &+ \int dx\, \langle\Psi, \Big[ c^*(\gamma_x) c(\gamma_x)+ c^*(\sigma_x) c(\sigma_x)+c^*(\gamma_x) c^*(\sigma_x)+c(\gamma_x) c(\sigma_x)\Big]\Psi\rangle. \end{split} \]
The claim now follows by (\ref{eq:choice}), combining Lemma \ref{lm:a-bds} and Lemma \ref{lm:c's}. 
\end{proof} 


\section{Energy of the trial state} 

In this section, we prove Prop. \ref{prop:Psi-energy}. To this end, we are going to estimate separately the kinetic and the potential energy of the trial state (\ref{eq:trial}). The bound for the kinetic energy was shown in  \cite[Prop. 2.5]{BBCOS}.
\begin{proposition}\label{prop:kin}
Let $\Psi$ be defined as in (\ref{eq:trial}) and $\cK$ denote the kinetic energy operator (\ref{eq:cK}). Then, we have 
\begin{equation}\label{eq:Kpsi} \begin{split} \frac{1}{L^3} \langle \Psi, \cK \Psi \rangle \leq \; &\rho_0^2 \int |\nabla f_\ell(x)|^2 dx  + 2\rho_0 \|\s\|^2  \int |\nabla f_\ell(x)|^2 dx  + \int |\nabla f_\ell (x)|^2 |\sigma (x)|^2 dx \\ &+ 2\int f_\ell (x) \nabla f_\ell (x) \sigma (x) \nabla \sigma (x) dx  +  \int |f_\ell(x)|^2 |\nabla \s(x)|^2 dx \\	& + 2\rho_0 \int |\nabla f_\ell(x)|^2 \big[(\g \ast \s)(x) + (\s \ast \s)(x)\big] dx \\
& + 2\rho_0 \int  f_\ell(x) \nabla f_\ell(x) \cdot \nabla \s (x)  dx\, + C \rho^{5/2+\delta}
\end{split}\end{equation} 
for $0 < \delta < \eps$ with $\eps > 0$ small enough.
\end{proposition}

In the next proposition, whose proof is deferred to Sect. \ref{sec:pot}, we bound the potential energy of (\ref{eq:trial}). This is the main novelty of this paper. 
\begin{proposition}\label{prop:pot}
Let $\Psi$ be defined as in (\ref{eq:trial}) and $\cV$ denote the potential energy operator (\ref{eq:cV}). Then, we have 
\begin{equation}\label{eq:Vpsi}\begin{split} 
\frac{1}{L^3} \langle \Psi, \cV \Psi \rangle \leq \; &\frac{\rho_0^2}{2} \int dx \, V(x) f_\ell^2 (x) + \rho_0 \| \sigma \|_2^2 \int dx \, V(x) f_\ell^2 (x) \\ &+ \rho_0 \int dx \, V(x) f_\ell^2 (x)  \big( (\sigma * \sigma) (x) + (\sigma * (\gamma-1)) (x) \big) + C \rho^{5/2+ \eps/2} 
\end{split} \end{equation}
for $\eps, \delta > 0$ small enough.
\end{proposition} 

With Prop. \ref{prop:kin} with Prop. \ref{prop:pot} we can now conclude the proof of Prop. \ref{prop:Psi-energy}. 
\begin{proof}[Proof of Prop. \ref{prop:Psi-energy}] Combining (\ref{eq:Kpsi}) and (\ref{eq:Vpsi}) and introducing the notation \begin{equation}
\cE_\ell(x) :=   |\nabla f_\ell(x)|^2 +  \frac 1 2 V(x)f^2_\ell(x)
\end{equation} 
we arrive at 
\begin{equation}\begin{split}
\frac{1}{L^3} \langle &\Psi, (\cK + \cV) \Psi \rangle \\ \leq\; & \rho_0^2 \int \cE_\ell(x)\, dx  
% \rho_0^2 \int |\nabla f_\ell(x)|^2 \dx  
+ 2\rho_0 \|\s\|^2  \int  \cE_\ell(x) dx  
%+ 2\rho_0 \|\s\|^2  \int |\nabla f_\ell(x)|^2 dx  
+ \int |\nabla f_\ell (x)|^2 |\sigma (x)|^2 dx \\ &+ 2\int f_\ell (x) \nabla f_\ell (x) \sigma (x) \nabla \sigma (x) dx  +  \int |f_\ell(x)|^2 |\nabla \s(x)|^2 dx \\
& + 2\rho_0 \int \cE_\ell(x) \big[(\g-\mathbbm{1}) \ast \s)(x) + (\s \ast \s)(x)\big] dx + 2\rho_0 \int |\nabla f_\ell(x)|^2 \sigma(x) dx \\
%& + 2\rho_0 \int |\nabla f_\ell(x)|^2 \big[(\g \ast \s)(x) + (\s \ast \s)(x)\big] dx \\	
& + 2\rho_0 \int  f_\ell(x) \nabla f_\ell(x) \cdot \nabla \s (x)  dx + C \rho^{5/2+ \delta} ,
\end{split}
	\end{equation}
if $0 < \delta \leq \eps /2$ and $\eps > 0$ is small enough. Defining 
\begin{equation}\label{def:nu}
\begin{split} 
\th (x) = &\; - \rho_0 \omega_\ell (x) + f_\ell (x) \tilde{\sigma} (x) = - \rho_0 \chi_\ell (2x)  \omega_\ell (x) + \chi_{\ell_0} (2x) (1 - \chi_\ell (2x)) s (x) 
%\\ = & \; \chi_{\ell_0}(2x) \Big[ -\chi_\ell(2x) \rho_0 \o_{\ell_0}(x) + (1- \chi_\ell(2x))s(x)\Big]
\end{split} 
\end{equation} 
with $\tilde{\sigma}$ as in (\ref{eq:tlsigma}), we obtain 
\begin{equation}\label{eq:Hpsi} \begin{split} 
\frac{1}{L^3} \langle &\Psi, (\cK + \cV) \Psi \rangle \\ \leq\; &\| \nabla \th\|_2^2 +\frac{\rho_0^2}{2} \int V(x)f^2_\ell(x) dx  + 16 \pi \aa \rho_0 \| s\|_2^2 + 8 \pi \aa \rho_0   ((g-\mathbbm{1}) \ast \sigma)(0) + \sum_{i=1}^4 \mathcal{R}_i 
\end{split}
\end{equation} 
with 
\[\begin{split}
\mathcal{R}_1 =\; &  2\rho_0 \|\s\|_2^2  \int \cE_\ell(x) dx  - 8 \pi \aa \rho_0 \| s\|_2^2 \\
\mathcal{R}_2 = \; & 2\rho_0 \int  \cE_\ell(x) ((\gamma-\mathbbm{1}) \ast \sigma)(x) dx - 8 \pi \aa \rho_0   ((g-\mathbbm{1}) \ast \sigma)(0) \\
\mathcal{R}_3 = \; & 2\rho_0 \int \cE_\ell(x)^2 (\sigma \ast \sigma)(x) dx - 8 \pi \aa \rho_0 \| s\|_2^2
\end{split}\]
and
\[
\mathcal{R}_4 =  \| \nabla \{- \rho_0 \o_\ell +f_\ell \s\}\|_2^2 - \| \nabla \{- \rho_0 \o_\ell +f_\ell \tilde\s\}\|_2^2\,.
\]
The bound $|\cR_4| \le C \rho^{5/2+\delta}$ was established in the proof of \cite[Lemma 3.1]{BBCOS}. The estimates $|\cR_i| \le C \rho^{5/2+\delta}, i=1,\ldots,3,$ follow from \eqref{eq:en-fl}, comparing $((\gamma-1) * \sigma) (x)$ with $((\gamma-1) * \sigma) (0)$ and then with $((g-1) * \sigma) (0)$, comparing $(\sigma * \sigma) (x)$ with $(\sigma * \sigma) (0)$ and  estimating 
\begin{equation}\label{eq:sigma-s}
\big| \| \sigma \|^2 - \| s \|^2 \big| \lesssim \rho^{3/2+\eps} \,.
\end{equation} 
The details can be found below \cite[Eq.~(3.6)]{BBCOS}, with  \eqref{eq:en-fl} replacing \cite[Eq.~(2.6)]{BBCOS}. 

Next, we observe that, with $\omega = 1 - f$ and $f$ the untruncated minimizer of (\ref{eq:Ef}), 
\begin{equation}\label{eq:nabla-nu}
\| \nabla \th\|^2 +\frac{\rho_0^2}{2} \int V(x)f^2_\ell(x) dx \leq  \|\nabla (\th + {\rho_{0}}\omega \chi_{\ell_0} (2\cdot)) \|^2 + 4\pi \frak{a} \rho^2_0 + C \rho^{5/2+\eps} \,.
\end{equation}
This follows from (see \cite[Proof of Lemma 3.2]{BBCOS} for more details) \[ |\langle \nabla (\th + \rho_0 \omega \chi_{\ell_0} (2 \cdot) ) , \nabla (\rho_0 \omega \chi_{\ell_0} (2 \cdot)) \rangle| \lesssim \rho^{5/2+\eps} \]
and from the estimate
\[ \begin{split}  
\rho_0^2 & \| \nabla (\omega \chi_{\ell_0} (2 \cdot)) \|^2  +\frac{\rho_0^2}{2} \int V(x)f^2_\ell(x) dx \\
&\leq \rho_0^2 \int \cE_{\ell_0}(x) dx + C \rho^{5/2+\eps} \leq 4\pi \frak{a} \rho_0^2 \big( 1 + C \frak{a} / \ell_0) + C \rho^{5/2+\eps} \leq 4\pi \frak{a} \rho_0^2 + C \rho^{5/2+\eps}\end{split}  \]
where we used (\ref{eq:en-fl}) and the fact that $f_\ell = f$, on $\text{supp } V$. From (\ref{eq:Hpsi}) we arrive therefore at  
\[ \begin{split} \frac{1}{L^3} \langle \Psi, (\cK + \cV) \Psi \rangle \leq \; &4\pi \frak{a} \rho_0^2 +  
\|\nabla (\th + {\rho_{0}}\omega \chi_{\ell_0} (2\cdot)) \|^2 \\ &+ 16 \pi \aa \rho_0 \| s\|_2^2 + 8 \pi \aa \rho_0   ((g-\mathbbm{1}) \ast \sigma)(0) + C \rho^{5/2+\eps} \,.
\end{split} \]
From this point, we proceed as below \cite[Eq. (3.11)]{BBCOS} to conclude that 
\[\frac{1}{L^3} \langle \Psi, (\cK + \cV) \Psi \rangle \leq 4 \pi \frak{a} \rho^2 \left( 1 + \frac{128}{15 \sqrt{\pi}} (\rho \frak{a}^3)^{1/2} \right) + \sum_{i=1}^5 \mathcal{E}_i + C \rho^{5/2+\eps}\,, \]
with the error terms 
\begin{align*}
\mathcal{E}_1 =\; & \left\|\nabla (\th + {\rho_{0}}\omega \chi_{\ell_0} (2\cdot)) \right\|^2 - \left\|\nabla (s + {\rho_{0}}\omega_{\ell_0}) \right\|^2\\
  \mathcal{E}_2 =\; & 8\pi  \mathfrak{a} {\rho_0} \|s\|^2 -\frac{1}{|\L|}\sum_{k\in \L^*_+} \rho_0\widehat{V}_\mathrm{eff}(k)  \widehat s_k^2 \\
 \mathcal{E}_3 =\; & 8\pi  \mathfrak{a} \rho_0\,   ((g-\mathbbm{1}) \ast s)(0) - \frac{1}{|\L|}\sum_{k\in \L^*_+} \rho_0\widehat{V}_\mathrm{eff}(k)\,  \widehat{(g-\mathbbm{1})}(k) \widehat s_k \\
\mathcal{E}_4 = \; &  \frac{\rho_{0}}{|\L|}\sum_{k\in \L^*_+} \big( 2|k|^2\widehat \omega_{\ell_0}(k) -\widehat{V}_\mathrm{eff}(k) \big) \widehat s_k \\
\mathcal{E}_5 =\; &  \frac{\rho_0^2}{|\L|}\sum_{k\in \L^*_+}\Big(  |k|^2  \widehat \omega^2_{\ell_0}(k) -  \frac{ \widehat{V}^2_\mathrm{eff}(k)}{4 |k|^2} \Big)\,.
\end{align*} To estimate $\mathcal{E}_1$ we first rewrite, as in \cite{BBCOS}, 
\begin{equation}\label{eq:E1-Z}
\th(x) + \rho_0 \omega (x) \chi_{\ell_0} (2x) = (1 -\chi_\ell (2x)) Z(x) - (1- \chi_{\ell_0} (2x)) Z(x)
\end{equation} 
with $Z(x)= \rho_0 \omega_{\ell_0} (x) +s(x)$.  Hence 
\begin{equation}  \label{Eq:Explicit_E_2} 
   \begin{split} 
       | \mathcal{E}_1| 
         \leq \; &\left\langle \nabla Z, \left[(1  -  \chi_\ell (2\cdot) )^2 -1 \right] \nabla Z \right\rangle  \\ &+ 2 \big|  \langle (1-\chi_\ell (2\cdot)) \nabla Z, -\frac{2}{\ell} \nabla \chi (2\cdot / \ell) Z(x) - \nabla \{ (1 - \chi_{\ell_0} (2\cdot)) Z \} \rangle\big|  \\ &+ C\ell^{-2}  \|  \nabla \chi (2\cdot / \ell) Z \|^2 + C \| \nabla \{ (1 - \chi_{\ell_0} (2\cdot)) Z \} \|^2 \,. \end{split}  \end{equation} 
To bound the r.h.s. of \eqref{Eq:Explicit_E_2} we decompose $Z = Z_1 + Z_2$, where $Z_1$ is defined through the Fourier coefficients
\be \label{def:Z1}
\widehat Z_1(k) =  \rho_0 \widehat{\o}_{\ell_0}(k)- \frac{\rho_0 \widehat{V}_{\rm{eff}}(k) }{2|k|^2}\,.   \ee
From Lemma \ref{lm:eta} and expanding (\ref{eq:hatsk}) for $|k| \gg \rho^{1/2}$, we find 
\[ |\widehat{Z}_2 (k)| \lesssim \frac{\rho}{k^2} \min \big\{ 1 , \frac{\rho}{k^2} \big\} \, . \]
This easily implies 
\[ \| Z_2 \|_\infty \lesssim \| \widehat{Z}_2 \|_1 \leq C \rho^{3/2} \]
and 
\[ \| \nabla Z_2 \|_p \lesssim \| k \widehat{Z}_2 \|_{p'} \lesssim \rho^{2-\frac{3}{2p}} \]
for all $3/2 < p < \infty$. On the other hand, recalling the definition (\ref{eq:Veff}) of $V_\text{eff}$, we find 
\[ Z_1 = \rho_0 \o_{\ell_0} - \rho_0 \omega = - \rho_0 \omega (1-\chi_{\ell_0}) \, . \]
Since $1- \chi_{\ell_0}$ is supported on $|x| \geq 2\ell_0$, where $\omega (x) = \frak{a}/|x|$, we obtain
\[ \| Z_1 \|_\infty \lesssim \rho / \ell_0 \lesssim \rho^{3/2} , \qquad \| \nabla Z_1 \|_p \lesssim \rho \ell_0^{-2+3/p} \]
for all $p > 3/2$. We conclude that \begin{equation}\label{eq:bds-Z} \| Z \|_\infty \lesssim \rho^{3/2}, \qquad \| \nabla Z \|_p \lesssim \rho^{2-3/(2p)} \end{equation}
for all $3/2 < p < \infty$ (in contrast with \cite{BBCOS}, we do not have a bound for the $L^\infty$-norm of $\nabla Z$, due to the weaker decay of $\widehat{Z}_2$ at infinity). With (\ref{eq:bds-Z}), we can estimate the r.h.s. of  (\ref{Eq:Explicit_E_2}). In particular, the first term is controlled by 
\[ \langle \nabla Z , [(1  -  \chi_\ell (2\cdot) )^2 -1 ] \nabla Z \rangle \lesssim \int_{|x| \leq 4\ell} |\nabla Z (x)|^2 dx \lesssim \ell^{3/p'} \| \nabla Z \|_{2p}^2  \lesssim  \ell^{3(1-1/p)} \rho^{4-3/(2p)} \]
for any $1 < p < \infty$. Choosing $p$ large enough, we find 
\[ \langle \nabla Z , [(1  -  \chi_\ell (2\cdot) )^2 -1 ] \nabla Z \rangle \lesssim  \rho^{5/2+\eps} \]
if $\eps,\delta > 0$ are small enough. The first term on the last line of (\ref{Eq:Explicit_E_2}) can be estimated by 
\[ \ell^{-2} \| \nabla \chi (2 \cdot / \ell) Z \|^2 \lesssim \ell^{-2} \int_{|x| \leq 4\ell} |Z (x) |^2 dx \lesssim \ell \| Z \|_\infty^2 \lesssim \ell \rho^3 \lesssim \rho^{5/2+\eps} \]
if $\eps > 0$ is small enough. As for the second term on the last line, we have 
\[ \begin{split} \| &\nabla (1- \chi_{\ell_0} (2 \cdot))Z \|^2 \\ &\lesssim \ell_0^{-2} \int_{\ell_0 \leq |x| \leq 4\ell_0} |Z (x)|^2 dx + \int_{|x| \geq \ell_0} |\nabla Z (x)|^2 \\ &\lesssim \rho_0^2 \ell_0^{-2} \int_{\ell_0 \leq |x| \leq 4\ell_0} \frac{\frak{a}^2}{|x|^2} dx + \ell_0^{-2} \int_{|x| \geq \ell_0} |s(x)|^2 dx + \rho_0^2 \int_{|x| \geq \ell_0} \frac{\frak{a}^2}{|x|^4} dx + \int_{|x| \geq \ell_0} |\nabla s (x)|^2\\ &\lesssim \rho^2 \ell_0^{-1} + \ell_0^{-4} \rho^{1/2} \lesssim \rho^{5/2+\eps} \end{split} \]
if $\eps, \delta > 0$ are small enough. In the last step, we used Lemma \ref{lm:eta} to bound the norms of $s, \nabla s$ restricted on $|x| \geq \ell_0$. The terms on the second line of (\ref{Eq:Explicit_E_2}) can be estimated similarly. We conclude that $|\cE_1| \lesssim \rho^{5/2+\eps}$, if $\eps ,\delta > 0$ are small enough. 

The estimates $|\mathcal{E}_i| \leq C \rho^{5/2+\d}$, for $i=2,\ldots,5$,  can be shown as below \cite[Eq. (3.24)]{BBCOS}. 



%***********************
%**********************
%
%
%proceed differently from \cite{BBCOS} and we %first give a closer look to the Fourier %coefficient  $\widehat{Z} (k)$ of $Z$. With %the notation 
%\be \label{def:Z1}
%	\widehat Z_1(k) =  \rho_0 \widehat{\o}
%_{\ell_0}(k)- \frac{\rho_0 \widehat{V}
%_{\rm{eff}}(k) }{2|k|^2}\,.
%    \ee
%we split $ \widehat Z(k) = \widehat Z_1(k) +
%\widehat Z_2(k)$. With Lemma \ref{lm:omega}, %we have   $|\widehat{Z}_i (k)| \leq C \rho / |%k|^2 $, for $i=1,2$. 
%Moreover, for $|k| \gg \rho^{1/2}$
%    \[
%    \big| \widehat  Z_2(k)\big| \leq  C \rho^2 %|k|^{-4}
%    \]
% which follows from the bound   
%       \[
%   \Big|\widehat{s}_k + \frac{\rho_0 \widehat %V_{\rm{eff}}(k)}{2|k|^2}\Big| \leq \frac{C 
%\rho_0^2( \widehat{V}_{\rm{eff}}(k))^2}{|k|%^4}\,.
%   \]
%   On the other hand, writing 
%\[ 2k^2 \widehat{\omega}_{\ell_0} (k) = 2 \int %\min \Big\{ 1 , \frac{\frak{a}}{|x|} \Big\} 
%\chi (x/\ell_0) \Delta e^{-ik\cdot x} dx \]
%and integrating by parts (with Gauss theorem), %we find 
%\begin{equation}\label{eq:ome-V} 2k^2 
%\widehat{\omega}_{\ell_0} (k) - \widehat{V}
%_\text{eff} (k) = \widehat{z}_1 (k \ell_0) 
%\end{equation} 
%with $z_1 (x) = 4 \frak{a} \nabla \chi (x) 
%\cdot x/|x|^3 - 2\frak{a} \Delta \chi (x) / 
%|x|$. This implies that $\widehat{Z}_1 (k) = 
%\rho_0 \widehat{z}_1 (k \ell_0) / 2k^2$.   
%
%   With these definitions we can bound, %proceeding as below \cite[Eq. (3.22)]{BBCOS}
%   \be \label{eq:Z-io}
%     \| Z \|_\infty \leq \| \widehat{Z} \|_1 
%\leq \| \widehat{Z}_1 \|_1 + \| \widehat{Z}_2 
%\|_1 \leq C \rho^{3/2}\,. \ee
%   We now set
%   \[
%   \hat Z_0(k):= \widehat{Z}_1(k) + \widehat{Z}
%_2(k) \mathds{1}(|k|\leq 1)\,, \qquad    \hat %Z_e(k):= \widehat{Z}_2(k) \mathds{1}(|k|> 1)
%   \]
%and we denote by $Z_0$ and $Z_e$ the functions %with Fourier coefficients $\hat Z_0$ and $\hat %Z_e$. It follows that
%\be \begin{split} \label{eq:nablaZ0-io} 
%\| \nabla Z_0\|_\infty &\leq \| |\cdot| 
%\widehat Z_0 \|_1 \leq \frac{C}{|\L|} 
%\sum_{\substack{k \in \L^*\\ |k| \leq %K\rho^{1/2}}} \frac{\rho}{|k|} +   \frac{C}{|
%\L|} \sum_{\substack{k \in \L^*\\ |k| \geq K\rho^{1/2}}}  \bigg[ \frac{\rho |\widehat z_1( k\ell_0 )|}{|k|}  + |k| |\widehat{Z}_2(k)| \mathds{1}(|k|\leq 1)  \bigg] \\ &\leq %C \rho^2 + C \rho \left[ \frac{1}{|\L|} \sum_{k\in \L^*} |\widehat{z}_1 (k \ell_0)|^2 |k|^2 \right]^{1/2} \left[ \frac{1}{|\L|} \sum_{|k| \geq K \rho^{1/2}}  \frac{1}{|k|^4} %\right]^{1/2}  \\ &\hspace{.4cm} + \frac{C \rho^2}{|\Lambda|} \sum_{K \rho^{1/2}\leq |k| \leq 1 } \frac{1}{|k|^3}  \\ &\leq C \rho^2 + C \rho^{3/4}  \ell_0^{-5/2}  + C 
%\rho^2 |\log \rho | \leq C \rho^2 |\log \rho |   
%\end{split}\ee
%where we used $\| |\cdot|^2\widehat{z}_1 (\ell_0 \cdot) \|_2 = \| \nabla z_1 (x/\ell_0) / \ell_0^3 \|_2 \leq C \ell_0^{-5/2}$.  On the other hand
%   \be \label{eq:nablaZe-2}
%   \| \nabla Z_e \|^2_2 =\sum_{\substack{k \in \L^*\\ |k| > 1 }} |k|^2 |\hat Z_2(k)|^2 \leq \frac{C}{|\L|} \sum_{\substack{k \in \L^*\\ |k| > 1}}  \frac{\rho^4}{|k|^6} \leq %C \rho^4 \,.
%   \ee
%With the splitting $Z=Z_0 + Z_e $ we can therefore bound the r.h.s. of \eqref{Eq:Explicit_E_2} as 
%%\[ 
%%|\left\langle \nabla Z, \left[(1  -  \chi_\ell (2\cdot) )^2 -1 \right] \nabla Z \right\rangle| \leq  C \| \nabla Z_e \|^2_2 + C \ell^3 \| \nabla Z_0 \|_\infty^2 \leq C 
%\rho^{4-3\d} 
%%\]
% \begin{equation}\label{eq:E1ZZ} 
% \begin{split}
% |\mathcal{E}_1|  \leq &\; C \ell^3 \| \nabla Z_0 \|_\infty^2  + C \| \nabla Z_e \|^2_2 + C \ell^2 \| \nabla Z_0 \|_\infty \| Z \|_\infty + \ell^{-1/2}\|\nabla Z_e\|_2 \| Z 
%\|_\infty\\
% & + C \ell \| Z \|_\infty^2 + C \| \nabla Z |_{|x| \geq \ell_0} \|^2 + C\ell_0^{-2} \| Z |_{|x| \geq \ell_0} \|^2 \,.
% \end{split}  \end{equation} 
% With \eqref{eq:Z-io}, \eqref{eq:nablaZ0-io}, \eqref{eq:nablaZe-2}, together with the bounds  
% \[ C \ell_0^{-2} \| Z |_{|x| \geq \ell_0} \|_2^2  \leq C\rho^{5/2+\eps} , \qquad C  \| \nabla Z |_{|x| \geq \ell_0} \|^2_2 \leq C \rho^{5/2+\eps} \]
% shown below \cite[Eq. (3.19)]{BBCOS}, we conclude that $|\mathcal{E}_1|\leq C \rho^{5/2+\eps}$.
%

\end{proof}


\section{Potential energy of the trial state}  
\label{sec:pot}

The goal of this section is to prove Proposition \ref{prop:pot}. To this end, we apply the identity (\ref{eq:id}) to write 
\begin{equation*}
\begin{split} 
a_y a_x \Psi =\; &\sqrt{\rho_0} a_y J(x) \Psi + a_y J(x) \int dz \, \nu (x-z) a_z^* J(z) \Psi \\ = \; &\sqrt{\rho_0} f_\ell(x-y) J(x) a_y \Psi + f_\ell(x-y) J(x) \int dz \, \nu (x-z) a_y a_z^* J(z) \Psi \\ = \; &\sqrt{\rho_0} f_\ell(x-y) J(x) a_y \Psi + f_\ell(x-y)  \nu (x-y)  J(x) J(y) \Psi \\ &+ f_\ell(x-y) J(x) \int dz \, \nu (x-z) f_\ell(z-y) a_z^* J(z) a_y \Psi \,.
\end{split} \end{equation*} 
Again with (\ref{eq:id}), we obtain 
\begin{equation}
\label{eq:axayPsi}
\begin{split}  
a_y &a_x \Psi \\ = \; & \rho_0 f_\ell(x-y) J(x) J(y) \Psi + f_\ell(x-y) \nu (x-y) J(x) J(y) \Psi \\ &+ \sqrt{\rho_0} f_\ell(x-y)  \int dz \, \nu (y-z) f_\ell(x-z) f_\ell(y-z) a_z^* J(x) J(y)  J(z) \Psi \\ &+ \sqrt{\rho_0} f_\ell(x-y) \int dz  \, \nu (x-z) f_\ell(y-z) f_\ell(x-z) a_z^* J(x) J(y) J(z)  \Psi \\ &+ f_\ell(x-y)  \int dz dw \, \nu (x-z) \nu (y-w) \\ & \hspace{.3cm} \times  f_\ell(x-z) f_\ell(y-z) f_\ell(x-w) f_\ell(y-w) f_\ell(z-w) a_z^* a_w^* J(x) J(y) J(z) J(w) \Psi 
\\ =: & \sum_{j=1}^5 \phi_j (x;y) \end{split} \end{equation} 
where we observe that $\phi_4 (x;y) = \phi_3 (y;x)$. 

In the following lemmas, we estimate the different contributions of the terms $\{ \phi_j \}_{j=1}^5$ to the expectation of the potential energy in the trial state $\Psi$. 
\begin{lemma} \label{lm:phi1} 
We have
\[ \frac{1}{2L^3} \int dx dy \, V(x-y) \| \phi_1 (x;y) \|^2 \leq \frac{\rho_0^2}{2} \int dx \, V(x) f_\ell^2 (x)\,.  \] 
\end{lemma} 
\begin{proof} 
Follows immediately from $0 \leq J(x), J(y) \leq 1$. 
\end{proof} 

\begin{lemma} \label{lm:phi2}
We have 
\[ \frac{1}{2L^3} \int dx dy \, V(x-y)  \| \phi_2 (x;y) \|^2 \leq C \rho^{3}. \] 
\end{lemma} 
\begin{proof} 
From $0 \leq J(x), J(y) \leq 1$, we immediately find 
\[  \frac{1}{2L^3} \int dx dy \, V(x-y)  \| \phi_2 (x;y) \|^2 \leq  \frac{1}{2} \int dx \, V(x) f_\ell^2 (x) |\nu (x)|^2  \,.\]
Next, we recall $\nu (x) = (\gamma^{-1} * \sigma) (x) = ((\gamma^{-1} -1) * \sigma) (x) + \sigma (x)$, where (from Lemma \ref{lm:eta}) 
\[ \| (\gamma^{-1} -1) * \sigma \|_\infty \leq \| \gamma^{-1} - 1 \|_2 \| \sigma \|_2 \lesssim \rho^{3/2} \]
and, on the support of $V$, 
\[ |\sigma (x)| = \ell_0^{-3} |\ph (x/ \ell_0)| \Big| \int \tilde{\sigma} (z) dz \Big| \lesssim \rho^{3/2+3\eps/2} \,.\]
We conclude that
\[ \frac{1}{2L^3} \int dx dy \, V(x-y)  \| \phi_2 (x;y) \|^2 \lesssim \rho^3\int V(x) f_\ell^2(x)dx \lesssim \rho^3. \]
\end{proof} 

\begin{lemma} \label{lm:phi3}
For $j=3,4$, we find 
\[ \frac{1}{2L^3} \int dx dy \, V(x-y) \| \phi_j (x;y) \|^2 \leq \frac{\rho_0}{2} \| \sigma \|_2^2 \int dx \, V(x) f_\ell^2 (x) + C \rho^{5/2+\eps} \]
if $\eps , \delta > 0$ are small enough. 
\end{lemma} 
\begin{proof} 
From the definition of $\phi_3$, we find  
%\[ \phi_3 (x;y) = \sqrt{\rho_0} f_\ell(x-y) \int dz \, \nu (y-z) f_\ell(x-z) f_\ell(y-z) a_z^* J(x) J(y) J(z) \Psi \]
\[ \begin{split}  \frac{1}{2} &\int dx dy \, V(x-y) \| \phi_3 (x;y) \|^2 \\ = \; & \frac{\rho_0}{2} \int dx dy dz dz' \, V(x-y) f_\ell^2 (x-y) \nu (y-z) \nu (y-z') \\ &\hspace{.3cm} \times f_\ell(x-z) f_\ell(y-z) f_\ell(x-z') f_\ell(y-z') \langle J(x) J(y) J(z) \Psi , a_z a_{z'}^* J(x) J(y) J(z') \Psi \rangle \,. \end{split}  \]
 Hence
 \[  \begin{split}  \frac{1}{2L^3} &\int dx dy \, V(x-y) \| \phi_3 (x;y) \|^2 \\ = \; &\frac{\rho_0}{2L^3} \int dx dy dz V(x-y) f_\ell^2 (x-y) \nu^2 (y-z) f_\ell^2 (x-z) f_\ell^2 (y-z) \| J(x) J(y) J(z) \Psi \|^2 \\ &+ \frac{\rho_0}{2L^3} \int dx dy dz dz' \, V(x-y) f_\ell^2 (x-y) \nu (y-z) \nu (y-z')  f_\ell^2 (z-z')  \\ &\hspace{.3cm} \times f_\ell^2(x-z) f_\ell^2(y-z) f_\ell^2 (x-z') f_\ell^2(y-z') \langle a_{z'} \Psi ,  J^2 (x) J^2 (y) J(z) J(z') a_z \Psi \rangle \\ =: \, & \text{A}_1 + \text{A}_2  \,.
 \end{split} \]
Estimating $0 \leq J (x) \leq 1$ and $0 \leq f_\ell(x) \leq 1$ for all $x \in \bR^3$, we obtain 
 \begin{equation}\label{eq:A1-fin} \text{A}_1 \leq \frac{\rho_0}{2} \| \nu \|^2 \int V(x) f_\ell^2 (x) dx \,. \end{equation} 
To bound $\text{A}_2$, we set $R = \rho^{-1/2-3\eps}$ and we observe that the contribution from $|y-z| > R$ is negligible, thanks to the fast decay of $\nu (y-z)$, from Lemma \ref{lm:eta} (in particular, Eq.~(\ref{eq:decay}); since $m \in \bN$ is arbitrary, we can gain any power of $\rho$). More details about this argument can be found in \cite{BBCOS}, following Eq. (5.11). To bound the contribution from $|y-z| \leq R$, we notice instead that
\[ \begin{split} \rho_0 \int_{|y-z| \leq R} &dx dy dz dz' \, V(x-y) f_\ell^2(x-y)| \nu (y-z')|^2  \big\langle a_z \Psi, \big( 1 - J^2 (x) J^2 (y) J(z) J(z') \big) a_z \Psi \big\rangle  \\ &\leq \rho_0 \int_{|y-z| \leq R} dx dy dz dz' dw \, V(x-y) f_\ell^2(x-y) |\nu (y-z')|^2 \\ & \hspace{2cm} \times ( u_\ell (x-w) + u_\ell (y-w) + \omega_\ell (z-w) + \omega_\ell (z' - w))  \| a_z a_w \Psi \|^2  \end{split} \]
where we used the notation $u_\ell = 1-f_\ell^2$ and $\omega_\ell = 1- f_\ell$. Since $u_\ell (s) = 0$ for $|s| > 4\ell$ and since $V$ has compact support, we can estimate the term proportional to $u_\ell (x-w)$ by 
\[ \rho_0  \int_{|z-w| \leq C R} dx dy dz dz' dw \, V(x-y) f_\ell^2(x-y)| \nu (y-z')|^2 u_\ell (x-w) \| a_z a_w \Psi \|^2 \lesssim \rho^{3-15\eps-2\delta} L^3 \]
where we used $\| u_\ell \|_1 \lesssim \ell^2$ and $\|Vf_\ell^2\|_1\lesssim 1$, (\ref{eq:zeta-combi1}) and Lemma \ref{lm:a-bds}. The contributions proportional to $u_\ell (y-w)$, $\omega_\ell (z'-w)$ can be bounded similarly (using also the fast decay of $\nu (y-z')$ to restrict the integral to $|y-z'| \leq R$, producing an additional negligible error term). As for the contribution proportional to $u_\ell (z-w)$, we use Lemma \ref{lm:aaaaC}. We find 
\[ \rho_0 \int_{|y-z| \leq R} dx dy dz dz' dw \,V(x-y)f_\ell^2(x-y) |\nu (y-z')|^2 u_\ell (z-w) \| a_z a_w \Psi \|^2 \lesssim \rho^{3-45\eps-\delta} L^3. \]
By Cauchy-Schwarz, we conclude that 
\[ \begin{split} \Big| \rho_0 &\int dx dy dz dz' \, V(x-y) f_\ell^2 (x-y) \nu (y-z) \nu (y-z')  f_\ell^2 (z-z')  f_\ell^2(x-z) f_\ell^2(y-z)  \\ &\hspace{2cm} \times f_\ell^2 (x-z') f_\ell^2(y-z') \langle a_{z'} \Psi ,  \big( 1- J^2 (x) J^2 (y) J(z) J(z')  \big) a_z \Psi \rangle \Big| \lesssim \rho^{5/2+\eps} L^3\end{split} \]
if $\eps, \delta > 0$ are small enough. Next, we observe that
 \begin{equation}\label{eq:phi31-f} 
 \begin{split} \rho_0 &\int dx dy dz dz' \, V(x-y) f_\ell^2(x-y)|\nu (y-z)| |\nu (y-z')| \\ &\hspace{1cm} \times \big( 1 - f_\ell^2 (z-z')  f_\ell^2(x-z) f_\ell^2(y-z) f_\ell^2 (x-z') f_\ell^2(y-z') \big) \| a_z \Psi \|^2 \\ &\lesssim 
 \rho_0 \int dx dy dz dz' \, V(x-y)f_\ell^2(x-y) |\nu (y-z)| |\nu (y-z')|  \\ &\hspace{1cm} \times \big( u_\ell (z-z') + u_\ell (x-z) + u_\ell (y-z) + u_\ell (x-z') + u_\ell (y-z') \big) \| a_z \Psi \|^2 .\end{split} \end{equation}
 The term proportional to $u_\ell (z-z')$ can be bounded by 
 \[\begin{split}
     \rho_0 \int dx &dy dz dz' \, V(x-y)f_\ell^2(x-y) |\nu (y-z)| |\nu (y-z')|  u_\ell (z-z') \| a_z \Psi \|^2\\
     \lesssim\;& \rho_0\|\nu\|_\infty \int dx dy dz dz' \, V(x-y)f_\ell^2(x-y) |\nu (y-z')|  u_\ell (z-z') \| a_z \Psi \|^2 
     \lesssim \rho^{3-6\eps- \delta} L^3
 \end{split}  \]
 using Lemma \ref{lm:a-bds}, $\|Vf_\ell^2\|_1\lesssim 1$,  and $\| \nu \|_\infty \lesssim \rho / \ell$, $\| \nu \|_1 \lesssim 1$ from (\ref{eq:zeta-combi1}). The other contributions can be bounded similarly. 
 %As for the term proportional to $u(x-z')$, we can estimate it by 
 %\[ \rho_0 \int dx dy dz dz' \, V(x-y) |\nu (y-z)| |\nu (y-z')|  u(x-z') \| a_z \Psi \|^2 \leq C \rho^{5/2} \ell^{-2} L^3 ,\] 
%since $\| \nu \|_\infty \lesssim \rho/\ell$ The term proportional to $u(y-z')$ on the r.h.s. of (\ref{eq:phi31-f}) can be bounded similarly. 
We conclude that 
 \[ \begin{split} \Big| &\rho_0 \int dx dy dz dz' \, V(x-y) f_\ell^2 (x-y) \nu (y-z) \nu (y-z') \\ & \times  \big(  f_\ell^2 (z-z')  f_\ell^2(x-z) f_\ell^2(y-z) f_\ell^2 (x-z') f_\ell^2(y-z') - 1 \big) \langle a_{z'} \Psi , a_z \Psi \rangle \Big| \lesssim \rho^{5/2+\eps} L^3\end{split} \]
if $\eps,\delta > 0$ are small enough. Hence, we arrive at 
\[ \text{A}_2 \leq C \rho^{3-\delta}  + \frac{\rho_0}{2L^3} \int dx dy dz dz' V(x-y) f_\ell^2 (x-y) \nu (y-z) \nu (y-z') \langle \Psi, a^*_{z'} a_z \Psi \rangle. \]
To estimate the last term, we write $a_z = \sqrt{\rho_0} + b_z, a_{z'} = \sqrt{\rho_0} +  b_{z'}$ and we observe that, since $\int \nu = 0$, the contribution of $\sqrt{\rho_0}$ vanishes. Therefore 
\[ \text{A}_2 \leq C \rho^{3-\delta}  + \frac{\rho_0}{2L^3} \int dx dy dz dz'\, V(x-y) f_\ell^2 (x-y) \nu (y-z) \nu (y-z') \langle \Psi, b^*_{z'} b_z \Psi \rangle .\]
Writing 
\[ b^*_{z'} = c^* (\gamma_{z'}) - c (\sigma_{z'}) , \quad b_z = c (\gamma_z) - c^* (\sigma_z) \, . \]
 we obtain 
 \[ \begin{split} \text{A}_2 \leq \; &C \rho^{3-\delta}  + \frac{\rho_0}{2L^3} \int dx dy dz dz' \,V(x-y) f_\ell^2 (x-y) \nu (y-z) \nu (y-z') \langle \sigma_{z'} , \sigma_z \rangle \\ &+ 
 \frac{\rho_0}{2L^3} \int dx dy dz dz' \,V(x-y) f_\ell^2 (x-y) \nu (y-z) \nu (y-z') \\ & \hspace{1cm} \times \langle \Psi, \big( c^* (\gamma_{z'}) c (\gamma_z) + c^* (\sigma_{z'}) c (\sigma_z) - c^* (\gamma_{z'}) c^* (\sigma_z) - c (\sigma_{z'}) c (\gamma_z) \big) \Psi \rangle .\end{split} \]
 With Cauchy-Schwarz and by Lemma \ref{lm:c's}, we find 
 \[ \rho_0 \int dx dy dz dz' \, V(x-y) f_\ell^2 (x-y) |\nu (y-z)| |\nu (y-z')| \| c (\tau_z) \Psi \|^2 \lesssim \rho^{4-25\eps-\delta} L^3 \]
 for both $\tau = \gamma, \sigma$ and 
 \[ \rho_0 \int dx dy dz dz' \, V(x-y) f_\ell^2 (x-y) |\nu (y-z)|^2 \| c (\gamma_{z'}) c (\sigma_z) \Psi \|^2 \lesssim \rho^{9/2-42\eps-\delta} L^3 .\]
Thus,
\[ \text{A}_2 \leq   \frac{\rho_0}{2} (\nu* \nu * \sigma * \sigma) (0) \int V(x) f_\ell^2 (x) dx + C \rho^{5/2+\eps}  \]
if $\eps , \delta> 0$ is small enough. To combine this term with (\ref{eq:A1-fin}), we observe that $\| \nu \|^2_2 = (\nu * \nu ) (0)$ and we recall that $\nu = \gamma^{-1} * \sigma$ to rewrite 
\[ (\nu * \nu) (0) + (\nu * \nu * \sigma * \sigma) (0) = (\nu * \nu * \gamma * \gamma) (0) = (\sigma * \sigma) (0) = \| \sigma \|_2^2 .\]
This concludes the proof of the lemma. 
 \end{proof} 

\begin{lemma} \label{lm:phi5} 
We have 
\[ \frac{1}{2L^3} \int dx dy \, V(x-y) \| \phi_5 (x;y) \|^2 \leq C \rho^{5/2+\eps} \]
if $\eps, \delta > 0$ are small enough.
\end{lemma} 
\begin{proof} 
From the definition of $\phi_5$, we obtain 
\[ \begin{split}  \int dx &dy \, V(x-y) \| \phi_5 (x;y) \|^2 \\ = \; & \int dx dy dz dw dz' dw' \, V(x-y) f_\ell^2 (x-y) \nu (x-z) \nu (y-w) \nu (x-z') \nu (y-w') \\ &\hspace{.3cm} \times f_\ell(x-z) f_\ell(x-z') f_\ell(y-z)  f_\ell(y-z') f_\ell(x-w)  f_\ell(x-w') f_\ell(y-w) f_\ell(y-w')  \\ &\hspace{.3cm} \times   f_\ell(z-w)f_\ell(z'-w') \langle J(x) J(y) J(z) J(w) \Psi, a_w a_z  a^*_{z'} a_{w'}^*  J(x) J(y) J(z') J(w') \Psi \rangle . \end{split} \] 
With the canonical commutation relations, we find 
\[ \begin{split} a_w a_z a_{z'}^* a^*_{w'} = \; &\delta (z-z') \delta (w-w') + \delta (w-z') \delta (z-w') + \delta (z-z') a_{w'}^* a_w + \delta (z'-w) a_{w'}^* a_z \\ &+ \delta (z-w') a_{z'}^* a_w + \delta (w-w') a_{z'}^* a_z + a_{z'}^* a_{w'}^* a_w a_z \,. \end{split} \]
Hence, we have 
\[ \frac{1}{2L^3} \int dx dy \, V(x-y) \| \phi_5 (x;y) \|^2 = \sum_{j=1}^5 \text{B}_j  \]
with 
\[ \begin{split} \text{B}_1 = \; &\frac{1}{2L^3} \int dx dy dz dw \, V(x-y) f_\ell^2 (x-y) \nu^2 (x-z)  \nu^2 (y-w)\\ &\hspace{.2cm} \times    f_\ell^2 (x-z) f_\ell^2 (x-w)  f_\ell^2 (y-z) f_\ell^2 (y-w)  f_\ell^2 (w-z) \| J(x) J(y) J(z) J(w) \Psi \|^2 \end{split} \]
\[ \begin{split} \text{B}_2 = \; &\frac{1}{2L^3} \int dx dy dz dw \,  V(x-y) f_\ell^2 (x-y) \nu (x-z) \nu (x-w) \nu (y-w) \nu (y-z) \\ &\hspace{.2cm} \times f_\ell^2 (x-z) f_\ell^2 (x-w)  f_\ell^2 (y-z) f_\ell^2 (y-w)  f_\ell^2 (w-z) \| J(x) J(y) J(z) J(w) \Psi \|^2 \end{split} \]
\[ \begin{split} \text{B}_3 =  \; &\frac{1}{L^3} \int dx dy dz dw dw' \, V(x-y) f_\ell^2 (x-y) \nu^2 (x-z) \nu (y-w) \nu (y-w') \\ &\hspace{.2cm} \times 
f_\ell^2(x-z) f_\ell^2 (y-z) f_\ell^2 (x-w) f_\ell^2 (x-w') f_\ell^2 (y-w) f_\ell^2 (y-w') f_\ell^2 (z-w) \\ &\hspace{.2cm} \times f_\ell^2 (z-w') f_\ell^2 (w-w') \langle a_{w'} \Psi, J(x)^2 J(y)^2 J(z)^2 J(w) J(w') a_w \Psi \rangle \end{split} \]
\[ \begin{split} \text{B}_4 = \; &\frac{1}{L^3}  \int dx dy dz dw dw' \, V(x-y) f_\ell^2 (x-y) \nu (x-z) \nu (x-w) \nu (y-w) \nu (y-w') \\ &\hspace{.2cm} \times 
f_\ell^2 (x-z)   f_\ell^2 (y-z) f_\ell^2 (x-w) f_\ell^2 (x-w') f_\ell^2 (y-w) f_\ell^2 (y-w') f_\ell^2 (z-w)  \\ &\hspace{.2cm} \times f_\ell^2 (z-w') f_\ell^2 (w-w') \langle a_{w'} \Psi, J(x)^2 J(y)^2 J(w)^2 J(z) J(w') a_z \Psi \rangle  \end{split} \]
\[ \begin{split} \text{B}_5 = \; &\frac{1}{2L^3} \int dx dy dz dw dz' dw' \, V(x-y) f_\ell^2 (x-y) \nu (x-z) \nu (x-z') \nu (y-w) \nu (y-w')  \\ &\hspace{.2cm} \times 
f_\ell^2 (x-z) f_\ell^2 (x-z') f_\ell^2 (x-w)  f_\ell^2 (x-w')  f_\ell^2 (y-z) ^2 (y-w) f_\ell^2 (y-w') \\ &\hspace{.2cm} \times f_\ell^2 (y-z') f_\ell^2 (z-w') f_\ell^2 (z-z') f_\ell^2 (z'-w) f_\ell^2 (w-w') f_\ell(z-w) f_\ell(z'-w')  \\ &\hspace{.2cm} \times \langle a_{z'} a_{w'} \Psi , J (x)^2 J(y)^2 J(z) J(w) J(z') J(w') a_z a_w \Psi \rangle .
 \end{split} \]
Since $0 \leq f_\ell(s) \leq 1, 0 \leq J(s) \leq 1$ for all $s \in \bR^3$, we immediately obtain 
\[ \text{B}_1 \leq \frac{\| \nu \|_2^4}{2} \int V(x) f_\ell^2 (x) dx \lesssim \rho^3 \]
and 
\[ \text{B}_2 \leq \frac{1}{2} \int V(x) f_\ell^2 (x) (\nu * \nu)^2 (x) dx \lesssim \| \nu * \nu \|_\infty^2 \int V(x) f_\ell^2(x)dx \lesssim \| \nu \|_2^4 \lesssim \rho^3.  \]

To bound the term $\text{B}_3$, we first observe that 
\[\begin{split}  \int &dx dy dz dw dw' V(x-y) f_\ell^2(x-y)\nu^2 (x-z) |\nu (y-w)| |\nu (y-w')|\\ &\hspace{2cm} \times \langle a_{w'} \Psi, \big( 1 - J^2 (x) J^2 (y) J^2 (z) J(w) J(w') \big) a_{w'} \Psi \rangle \\ \leq \; &\int dx dy dz dw dw'ds \, V(x-y) f_\ell^2(x-y)\nu^2 (x-z) |\nu (y-w)| | \nu (y-w')| \\ &\hspace{2cm} \times \big( u_\ell (x-s) + u_\ell (y-s) + u_\ell (z-s) + u_\ell (w-s) + u_\ell (w' -s) \big) \| a_{w'} a_s \Psi \|^2. \end{split} \]
The contribution proportional to $u_\ell (x-s)$ can be estimated with (\ref{eq:zeta-combi1}) and Lemma \ref{lm:a-bds}. Using also the decay (\ref{eq:decay}) of $\nu$ to restrict the integral to $|s-w'| \leq R = \rho^{-1/2-3\eps}$,  we find  
\[ \begin{split}  \int dx &dy dz dw dw' ds \, V(x-y)f_\ell^2(x-y) \nu^2 (x-z) |\nu (y-w)| | \nu (y-w')| u_\ell (x-s) \| a_{w'} a_s \Psi \|^2 \\  \leq \; & \| \nu \|_\infty \| \nu \|_2^2 \| \nu \|_1 \| u_\ell \|_1 \|Vf_\ell^2\|_1\int_{|s-w'| \leq R} dw' ds \, \| a_{w'} a_s \Psi \|^2 + C \rho^3 L^3 \lesssim C \rho^{3-18\eps-\delta} L^3.  \end{split} \]
The contributions proportional to $u_\ell (y-s), u_\ell (z-s), u_\ell (w-s)$ can be controlled analogously. To bound the term proportional to $u_\ell (w'-s)$, we apply Lemma \ref{lm:aaaaC}. We find 
\[ \begin{split}  \int dx dy dz &dw dw' ds \, V(x-y)f_\ell^2(x-y) \nu^2 (x-z) |\nu (y-w)| | \nu (y-w')| u_\ell (w'-s) \| a_{w'} a_s \Psi \|^2 \\  \leq \; & \| \nu \|^2_1 \| \nu \|_2^2 \|Vf_\ell^2\|_1\int_{|w'-s| \leq \ell} dw'ds \, \| a_{w'} a_s \Psi \|^2 \lesssim \rho^{7/2-42\eps-\delta} L^3. \end{split} \]
By Cauchy-Schwarz, we conclude that 
\[ \begin{split} \text{B}_3 \leq \; &\frac{1}{L^3} \int  dx dy dz dw dw' \, V(x-y) f_\ell^2 (x-y) \nu^2 (x-z) \nu (y-w) \nu (y-w') \\ &\hspace{.2cm} \times 
f_\ell^2(x-z) f_\ell^2 (y-z) f_\ell^2 (x-w) f_\ell^2 (x-w') f_\ell^2 (y-w) f_\ell^2 (y-w') f_\ell^2 (z-w) \\ &\hspace{.2cm} \times f_\ell^2 (z-w') f_\ell^2 (w-w') \langle a_{w'} \Psi,  a_w \Psi \rangle  + C \rho^{5/2+\eps} \end{split} \]
if $\eps, \delta > 0$ are small enough. Next, we estimate 
\[ \begin{split}  1  - \prod_{r,s} f_\ell^2 (r-s) \leq \; &u_\ell (x-z) + u_\ell (y-z) + u_\ell (x-w) + u_\ell (x-w') +u_\ell (y-w) \\ &+ u_\ell (y-w') + u_\ell (z-w) + u_\ell (z-w') + u_\ell (w-w') \end{split} \]
where the product runs over all $r, s \in \{ x,y,z,w,w'\}$, with $r \not = s$ and $(r,s) \not = (x,y)$. Considering for example the contribution proportional to $u_\ell (x-z)$, we can bound it by 
\begin{equation} \label{eq:uxz} \begin{split}  \int dx &dy dz dw dw' \, V(x-y)f_\ell^2(x-y) \nu^2 (x-z) |\nu (y-w)| |\nu (y-w')| u_\ell (x-z) \| a_{w'} \Psi \|^2 \\  &\leq \| \nu \|_\infty^2 \int dx dy dz dw dw' \, V(x-y) f_\ell^2(x-y)|\nu (y-w)| |\nu (y-w')| u_\ell (x-z) \| a_{w'} \Psi \|^2 \\  &\leq \| \nu \|_\infty^2 \| \nu \|_1^2 \| u_\ell \|_1\|Vf_\ell^2\|_1 \int dw' \, \| a_{w'} \Psi \|^2 \leq  \rho^{3-9\eps}  L^3 \end{split} \end{equation}
with (\ref{eq:zeta-combi1}) and Lemma \ref{lm:a-bds}. As for the term proportional to $u_\ell (x-w)$, we estimate it by 
\begin{equation}\label{eq:uxw} \begin{split}  \int dx dy &dz dw dw' \, V(x-y) f_\ell^2(x-y) \nu^2 (x-z) |\nu (y-w)|\\
&\qquad\qquad\qquad\times|\nu (y-w')| u_\ell (x-w) \| a_w' \Psi \|^2 \\ 
\leq\;& \| \nu \|_\infty   \int dx dy dz dw dw' \, V(x-y)f_\ell^2(x-y) \nu^2 (x-z)\\
&\qquad\qquad\qquad\times|\nu (y-w')| u_\ell (x-w) \| a_w' \Psi \|^2 \\ \leq\;& \| \nu \|_\infty \| \nu \|_2^2 \| \nu \|_1 \| u_\ell \|_1\|Vf_\ell^2\|_1 \int dw' \, \| a_{w'} \Psi \|^2 \leq  \rho^{7/2-6\eps-\delta}  L^3. \end{split} \end{equation}
The other contributions can be bounded either similarly to (\ref{eq:uxz}) or to (\ref{eq:uxw}). Hence, 
\[\begin{split}
    \text{B}_3 \leq \;&\frac{1}{L^3} \int dx dy dz dw dw' \, V(x-y) f_\ell^2 (x-y)  \nu^2 (x-z) \nu (y-w) \nu (y-w')  \langle a_{w'} \Psi,  a_w \Psi \rangle  
\\
&+ C \rho^{5/2+\eps},
\end{split}  \]
if $\eps, \delta > 0$ are small enough. We write $a_w = \sqrt{\rho_0} + b_w, a_{w'} = \sqrt{\rho_0} + b_{w'}$ and we observe that the contribution of $\sqrt{\rho_0}$ vanishes, since $\int \nu = 0$. Therefore, with Lemma \ref{lm:a-bds} and (\ref{eq:zeta-combi1}), we obtain 
\[ \begin{split} \text{B}_3 \leq \; &C\rho^{5/2+\eps} + \frac{1}{L^3} \int  dx dy dz dw dw' \, V(x-y) f_\ell^2(x-y) \nu^2 (x-z) |\nu (y-w)| \\
&\qquad\qquad\qquad\qquad\qquad\qquad\qquad\times|\nu (y-w')| \| b_{w'} \Psi \|^2 \\
 \leq \; & C \rho^{5/2+\eps} + \frac{C \| \nu \|_2^2 \| \nu \|_1^2\|Vf_\ell^2\|_1}{L^3} \int dw' \| b_{w'} \Psi \|^2 \leq C  \rho^{5/2+\eps} \end{split}  \] 
if $\eps, \delta > 0$ are small enough. 

To estimate $\text{B}_4$, we proceed similarly to what we did in order to remove from $\text{B}_3$ the factor $J^2 (x) J^2 (y) J^2 (w) J(z) J(w')$ and all factors of $f_\ell^2$ except for $f_\ell^2(x-y)$. We arrive at 
\[  \begin{split} \text{B}_4 \leq \; &C \rho^{5/2+\eps} + \frac{1}{L^3} \int dx dy dz dw dw' \, V (x-y) f_\ell^2 (x-y)  \nu (x-z) \nu (x-w) \\ &\hspace{6cm} \times \nu (y-w) \nu (y-w') \langle a_{w'} \Psi , a_z \Psi \rangle \, . \end{split} \]
As above, in the integral we can replace $a_{w'}, a_z$ by $b_{w'}, b_z$. With (\ref{eq:zeta-combi1}) and Lemma \ref{lm:a-bds}, we conclude that 
\[ \begin{split} \text{B}_4 \leq \; &C \rho^{5/2+\eps}+\\
&+\frac{1}{L^3} \int dx dy dz dw dw' \, V(x-y)f_\ell^2(x-y) |\nu (x-w)|^2 |\nu (x-z)| |\nu (y-w')| \| b_{w'} \Psi \|^2 \\ &+ \frac{1}{L^3}  \int dx dy dz dw dw' \, V(x-y)f_\ell^2(x-y) |\nu (y-w)|^2 |\nu (x-z)| |\nu (y-w')| \| b_{z} \Psi \|^2  \\ \lesssim \; &\frac{\| \nu \|_2^2 \| \nu \|_1^2 \|Vf_\ell^2\|_1}{L^3}  \int dw' \, \| b_{w'} \Psi \|^2 + C \rho^{5/2 + \eps} \leq C \rho^{5/2+\eps} \end{split} \]
if $\eps, \delta > 0$ are small enough. 

Finally, we consider the term $\text{B}_5$. As above, we first remove the $J$-operators and the $f_\ell^2$-factors $\prod_{r,s} f_\ell^2 (r-s)$. We find 
\[ \begin{split} \text{B}_5 \leq \; &C \rho^{5/2+\eps} + \frac{1}{L^3} \int dx dy dz dw dz' dw' \, V(x-y)f_\ell^2(x-y)  \nu (x-z) \nu (x-z') \nu (y-w)  \\ &\hspace{8cm} \times \nu (y-w')\langle a_{z' }a_{w'} \Psi, a_z a_w \Psi \rangle \\ = \; &C \rho^{5/2+\eps}+ \frac{1}{L^3} \int dx dy dz dw dz' dw' \, V(x-y) f_\ell^2(x-y) \nu (x-z) \nu (x-z') \nu (y-w) \\ &\hspace{8cm} \times \nu (y-w') \langle b_{z' }b_{w'} \Psi ,b_z b_w \Psi \rangle \,. \end{split} \]
Using the fast decay (\ref{eq:decay}) of $\nu$ to restrict the integrals to $|x-z| \leq R$, $|y-w| \leq R$, Cauchy-Schwarz, (\ref{eq:zeta-combi1}) and Lemma \ref{lm:a-bds}, we obtain 
\[ \begin{split} \text{B}_5 \leq \; &C \rho^{5/2+\eps} +\frac{C}{L^3} \int_{|x-z| \leq R, |y-w| \leq R} dx dy dz dw dz' dw' \, V(x-y) f_\ell^2(x-y)\\ &\hspace{7cm} \times |\nu (x-z')|^2 |\nu (y-w')|^2  \| b_z b_w \Psi \|^2\\  \leq \; &C \rho^{5/2+\eps} + C  \frac{R^3 \| \nu \|^4_2 \|Vf_\ell^2\|_1}{L^3} \int_{|z-w| \leq 2R} dz dw \, | b_z b_w \Psi \|^2 \leq C \rho^{5/2+\eps} 
  \end{split} \]
if $\eps > 0$ is small enough. This concludes the proof of the lemma. 
\end{proof} 

%off diagonal contributions

\begin{lemma} We have 
\[ \frac{1}{L^3} \text{Re } \int dx dy \, V(x-y) \langle \phi_1 (x;y) , \phi_2 (x;y) \rangle \leq \rho_0 \int dx \, V(x) f_\ell^2 (x) \nu (x)  + C \rho^{5/2+\eps}  \]
if $\eps , \delta > 0$ are small enough. 
\end{lemma} 
\begin{proof} 
From the definition of $\phi_1, \phi_2$, we obtain 
\[ \int dx dy \, V(x-y) \langle \phi_1 (x;y) , \phi_2 (x;y) \rangle = \rho_0 \int dx dy \, V(x-y) f_\ell^2 (x-y) \nu (x-y) \| J(x) J(y) \Psi \|^2 . \]
With (\ref{eq:zeta-combi1}), Lemma \ref{lm:a-bds} and $\| u_\ell \|_1 \lesssim \rho^{-2\delta}$, we find 
\[ \begin{split} \Big| \rho_0 \int dx dy \,& V(x-y) f_\ell^2 (x-y) \nu (x-y) \langle \Psi, \big( 1 - J^2 (x) J^2 (y) \big) \Psi \rangle \Big| \\ &\leq \rho_0 \int dx dy ds \, V(x-y) f_\ell^2(x-y)|\nu (x-y)| \big(u_\ell (x-s) + u_\ell (y-s) \big) \| a_s \Psi \|^2 \\ &\leq C \rho_0 \| \nu \|_\infty \| u_\ell \|_1 \|Vf_\ell^2\|_1 \int ds \, \| a_s \Psi \|^2  \leq C \rho^{5/2+\eps} L^3 \end{split} \]
if $\eps , \delta > 0$ are small enough. Hence,
\[ \frac{1}{L^3} \text{Re} \int dx dy \, V(x-y) \langle \phi_1 (x;y) , \phi_2 (x;y) \rangle \leq \rho_0 \int dx \, V(x) f_\ell^2 (x) \nu (x) + C \rho^{5/2+\eps} .\]
\end{proof}



\begin{lemma}\label{lm:phi1phi3} 
For $j=3,4$, we have 
\[ \frac{1}{L^3} \text{Re } \int dx dy \, V(x-y) \langle \phi_1 (x;y) , \phi_j (x;y) \rangle \leq C \rho^{5/2+\eps} \]
if $\eps,\delta > 0$ are small enough.
\end{lemma} 
\begin{proof} 
We consider only $j=3$, the case $j=4$ can be handled analogously. From the definition of $\phi_1, \phi_3$, we have
\[ \begin{split}  &\int dx dy \, V(x-y) \langle \phi_1 (x) , \phi_3 (x) \rangle \\ &= \rho_0^{3/2} \int dx dy dz \, V(x-y) f_\ell^2 (x-y) \nu (y-z) f_\ell^2 (y-z) f_\ell^2 (x-z) \langle a_z \Psi , J^2 (x) J^2 (y) J(z) \Psi \rangle .\end{split} \]
To estimate this contribution, we observe that, applying (\ref{eq:decay}) to restrict the integral to $|y-z| \leq R$ and then using Cauchy-Schwarz, 
\[ \begin{split} 
\Big|  \rho_0^{3/2} &\int dx dy dz \, V(x-y) f_\ell^2 (x-y) \nu (y-z) f_\ell^2 (y-z) f_\ell^2 (x-z) \\ &\hspace{7cm} \times  \langle a_z \Psi , \big(1 - J^2 (x) J^2 (y) J(z)\big)  \Psi \rangle \Big| \\  \leq \; &C \rho^{5/2+\eps} + \rho^{3/2} \Big( \int_{|y-z| \leq R} dx dy dz \, V(x-y)f_\ell^2(x-y) \\
&\hspace{5cm}\times\big( u_\ell (x-s) + u_\ell (y-s) + u_\ell (z-s) \big) \| a_z a_s \Psi \|^2 \Big)^{1/2} \\ &\hspace{.2cm} \times \Big( \int dx dy dz \, V(x-y)f_\ell^2(x-y) |\nu (y-z)|^2 \\
&\hspace{4cm}\times\big( u_\ell (x-s) + u_\ell (y-s) + u_\ell (z-s) \big) \| a_s \Psi \|^2 \Big)^{1/2} \\ \leq \; &C \rho^{5/2+\eps} + C \rho^{3/2} \|Vf_\ell^2\|_1 \bigg( \| u_\ell \|_1 \int_{|z-s| \leq R} dz ds \, \| a_z a_s \Psi \|^2\\
&\hspace{3.5cm}+ R^3 \int_{|z-s| \leq \ell} dz ds \, \| a_z a_s \Psi \|^2 \bigg)^{1/2} \bigg( \| \nu \|_2^2 \| u_\ell \|_1 \int ds \| a_s \Psi \|^2 \bigg)^{1/2} \\ \leq \; &C \rho^{5/2+\eps} 
\end{split} \]
if $\eps,\delta > 0$ are small enough, where, in the last step, we used (\ref{eq:zeta-combi1}), Lemma \ref{lm:a-bds} and Lemma~\ref{lm:aaaaC}. Furthermore, we notice that
\[ \begin{split} 
\Big|  \rho_0^{3/2} &\int dx dy dz \, V(x-y) f_\ell^2 (x-y) \nu (y-z) \big( 1 - f_\ell^2 (y-z) f_\ell^2 (x-z) \big) \langle a_z \Psi ,  \Psi \rangle \Big| \\ &\leq  \rho^{3/2} \int dx dy dz \, V(x-y)f_\ell^2(x-y)  |\nu (y-z)| \big( u_\ell (y-z) + u_\ell (x-z) \big) \| a_z \Psi \|  \\ &\leq \rho^{3/2} \| \nu \|_\infty \| u_\ell \|_1 \|Vf_\ell^2\|_1\int dz \| a_z \Psi \| \leq \rho^{5/2-\delta} L^{3/2} \Big( \int dz \, \|a_z \Psi \|^2 \Big)^{1/2} \leq \rho^{5/2+\eps} L^3 \end{split} \]
if $\eps, \delta > 0$ are small enough. The claim now follows from the observation that 
\[ \int dx dy dz \, V(x-y) f_\ell^2 (x-y) \nu (y-z) \langle a_z \Psi ,  \Psi \rangle = 0 \]
since $\int \nu = 0$. 
\end{proof}



\begin{lemma} \label{lm:phi1phi5} 
We have 
\[ \frac{1}{L^3} \text{Re } \int dx dy \, V(x-y) \langle \phi_1 (x;y) , \phi_5 (x;y) \rangle \leq \rho_0 \int dx \, V(x) f_\ell^2 (x) (\nu * \sigma * \sigma) (x) +C \rho^{5/2+\eps} \]
if $\eps, \delta > 0$ are small enough. 
\end{lemma} 
\begin{proof} 
From the definition of $\phi_1, \phi_5$, we have 
\[ \begin{split} 
 \int &dx dy \, V(x-y) \langle \phi_1 (x;y) , \phi_5 (x;y) \rangle \\ &= \rho_0 \int dx dy dz dw \, V(x-y) f_\ell^2 (x-y) \nu (x-z) \nu (y-w) \\ &\hspace{.5cm} \times  f_\ell^2 (x-z) f_\ell^2 (x-w) f_\ell^2 (y-z) f_\ell^2 (y-w) f_\ell(z-w) \langle a_z a_w \Psi, J^2 (x) J^2 (y) J(z) J(w) \Psi \rangle .\end{split} \]
 To bound this contribution, we first observe that 
 \[ \begin{split} 
 \Big| \rho_0 \int &dx dy dz dw \, V(x-y) f_\ell^2 (x-y) \nu (x-z) \nu (y-w)  f_\ell^2 (x-z) f_\ell^2 (x-w)  \\ &\hspace{.5cm} \times f_\ell^2 (y-z) f_\ell^2 (y-w) f_\ell(z-w) \langle a_z a_w \Psi, \big(1 - J^2 (x) J^2 (y) J(z) J(w) \big) \Psi \rangle \Big| \\
 \leq \; &C \rho^{5/2+\eps} L^3 + \rho_0 \Big( \int_{|y-w| , |x-z| \leq R} dx dy dz dw ds \, V(x-y)f_\ell^2(x-y) \, \\ &\hspace{3cm} \times \big(u_\ell (x-s) + u_\ell (y-s) + u_\ell (z-s) + u_\ell (w-s) \big) \| a_s a_w a_z \Psi \|^2 \Big)^{1/2} \\ &\times \Big(   \int_{|y-w| , |x-z| \leq R} dx dy dz dw ds \, V(x-y) f_\ell^2(x-y)\, |\nu (x-z)|^2 |\nu (y-w)|^2 \\ &\hspace{3cm} \times \big(u_\ell (x-s) + u_\ell (y-s) + u_\ell (z-s) + u_\ell (w-s) \big) \| a_s \Psi \|^2 \Big)^{1/2}  \\ \leq \; &C \rho^{5/2+\eps} L^3 + C \rho_0 \|Vf_\ell^2\|_1 \Big( \| u_\ell \|_1 \int_{|s-w|, |s-z| \leq C R}  ds dw dz \, \| a_s a_w a_z \Psi \|^2 \\ &\hspace{1cm} + R^3 \int_{|s-z| \leq \ell, |s-w| \leq CR} ds dw dz \| a_s a_w a_z \Psi \|^2 \Big)^{1/2} \Big( \| \nu \|_2^4 \| u_\ell \|_1 \int ds \, \| a_s \Psi \|^2 \Big)^{1/2} \\ \leq \; &C \rho^{5/2+\eps} L^3  
 \end{split} \]
if $\eps , \delta > 0$ are small enough. Here we used (\ref{eq:zeta-combi1}), Lemma \ref{lm:a-bds} and Lemma \ref{lm:aaaaC}. We can also estimate 
\[ \begin{split} 
 \Big| \rho_0 \int &dx dy dz dw \, V(x-y) f_\ell^2 (x-y) \nu (x-z) \nu (y-w) \\ &\hspace{.5cm}\times \big( 1-  f_\ell^2 (x-z) f_\ell^2 (x-w) f_\ell^2 (y-z) f_\ell^2 (y-w) f_\ell(z-w)\big)  \langle a_z a_w \Psi, \Psi \rangle \Big| \\ \leq \; &C \rho^{5/2+\eps} L^3+ \rho_0 \Big( \int_{|y-w| , |x-z| \leq R} dx dy dz dw \, V(x-y)f_\ell^2(x-y) \, \\ &\hspace{1cm} \times \big(u_\ell (x-z) + u_\ell (x-w) + u_\ell (y-z) + u_\ell (y-w) +u_\ell (z-w) \big) \| a_w a_z \Psi \|^2 \Big)^{1/2} \\ &\times \Big(   \int_{|y-w| , |x-z| \leq R} dx dy dz dw \, V(x-y)f_\ell^2(x-y) \, |\nu (x-z)|^2 |\nu (y-w)|^2 \\ &\hspace{1cm} \times  \big(u_\ell (x-z) + u_\ell (x-w) + u_\ell (y-z) + u_\ell (y-w) +u_\ell (z-w) \big) \Big)^{1/2}  \\ \leq \; &C \rho^{5/2+\eps}L^3 + C \rho_0 \|Vf_\ell^2\|_1\Big( \| u_\ell \|_1 \int_{|w-z| \leq C R}  dw dz \, \| a_w a_z \Psi \|^2 \\
 &\hspace{4cm}+ R^3 \int_{|w-z| \leq \ell} dw dz \| a_w a_z \Psi \|^2 \Big)^{1/2}  \Big( \| \nu \|_\infty^2 \| \nu \|_2^2 \| u_\ell \|_1 L^3 \Big)^{1/2} \\ \leq \; &C \rho^{5/2+\eps} L^3
 \end{split} \]
 if $\eps, \delta > 0$ are small enough. It follows that 
 \[ \begin{split}  \text{Re } &\int dx dy \, V(x-y) \langle \phi_1 (x;y) , \phi_5 (x;y) \rangle \\ &\leq \text{Re } \rho_0  \int dx dy dz dw \, V(x-y) f_\ell^2 (x-y) \nu (x-z) \nu (y-w) \langle a_z a_w \Psi, \Psi \rangle + C \rho^{5/2+\eps} L^3  \\ &= \text{Re } \rho_0 \int dx dy dz dw \, V(x-y) f_\ell^2 (x-y) \nu (x-z) \nu (y-w) \langle b_z b_w \Psi, \Psi \rangle + C \rho^{5/2+\eps} L^3 \end{split} \]
 because $\int \nu = 0$. Next, we write $b_s = c(\gamma_s) - c^* (\sigma_s)$, for $s = z,w$. We conclude that 
 \begin{equation}\label{eq:IVa} \begin{split}  \text{Re } &\int dx dy \, V(x-y) \langle \phi_1 (x;y) , \phi_5 (x;y) \rangle \\ &\leq  C \rho^{5/2+\eps} L^3 + \text{Re } \rho_0  \int dx dy dz dw \, V(x-y) f_\ell^2 (x-y) \nu (x-z) \nu (y-w)  \langle \gamma_z ; \sigma_w \rangle \\ &+\text{Re } \rho_0  \int dx dy dz dw \, V(x-y) f_\ell^2 (x-y) \nu (x-z) \nu (y-w) \\ &\hspace{2cm} \times \langle \big( c^* (\sigma_z) c (\gamma_w) + c^* (\sigma_w) c (\gamma_z) + c^* (\sigma_z) c^* (\sigma_w) + c(\gamma_z) c (\gamma_w)\big) \Psi, \Psi \rangle . \end{split} \end{equation} 
With Lemma \ref{lm:c's}, we remark that 
\[ \begin{split} 
\Big| &\rho_0 \int dx dy dz dw \, V(x-y) f_\ell^2 (x-y) \nu (x-z) \nu (y-w) \langle c (\gamma_w) \Psi , c (\sigma_z) \Psi \rangle \Big| \\ &\leq \rho_0 \int dx dy dz dw \, V(x-y)f_\ell^2(x-y) |\nu (x-z)| |\nu (y-w)| \left[ \| c (\gamma_w) \Psi \|^2 + \| c (\sigma_z) \Psi \|^2 \right] \\ &\leq C \rho \| \nu \|_1^2 \|Vf_\ell^2\|_1 \sum_{\tau = \gamma, \sigma} \int dw \, \| c (\tau_w) \Psi \|^2 \leq C \rho^{5/2+\eps} L^3 \end{split} \]
if $\eps , \delta > 0$ are small enough. Furthermore,
\[ \begin{split} 
&\Big| \rho_0 \int dx dy dz dw \, V(x-y) f_\ell^2 (x-y) \nu (x-z) \nu (y-w) \langle c (\gamma_w) c (\gamma_z) \Psi , \Psi \rangle \Big| \\ &\leq C \rho^{5/2+\eps} L^3 + \rho_0 \Big( \int_{|w-y| \leq R}  dx dy dz dw \, V(x-y) f_\ell^2(x-y) |\nu (x-z)|  \| c (\gamma_w) c (\gamma_z) \Psi \|^2 \Big)^{1/2}\\ &\hspace{3cm} \times  \Big(  \int dx dy dz dw \, V(x-y) f_\ell^2(x-y) |\nu (x-z)| |\nu (y-w)|^2 \Big)^{1/2}  \\ &\leq C \rho^{5/2+\eps} L^3 + C \rho_0 \|Vf_\ell^2\|_1 \Big(\| \nu \|_1 \int_{|w-z| \leq R} dw dz \, \| c (\gamma_w) c (\gamma_z) \Psi \|^2 \Big)^{1/2} (\| \nu \|_1 \| \nu \|_2^2 L^3)^{1/2} \\ &\leq C \rho^{5/2+\eps} L^3 . \end{split} \]
From \eqref{eq:IVa}, we find that 
\[ \begin{split} 
 \text{Re } \int dx dy \, V(x-y) \langle &\phi_1 (x;y) , \phi_5 (x;y) \rangle \\ &\leq  C \rho^{5/2+\eps} L^3 + \rho_0 L^3  \int dx \, V(x) f_\ell^2 (x) (\nu * \gamma * \sigma * \nu) (x)  
  \end{split} \] 
  which concludes the proof of the lemma, since $\nu * \nu * \gamma * \sigma = \nu * \sigma * \sigma$. 
\end{proof}

\begin{lemma} For $j=3,4$, we have 
\[ \frac{1}{L^3} \text{Re } \int dx dy \, V(x-y) \langle \phi_2 (x;y) , \phi_j (x;y) \rangle \leq C \rho^{5/2+\eps} \]
if $\eps ,\delta > 0$ are small enough. 
\end{lemma} 
\begin{proof}
The proof is very similar to the proof of Lemma \ref{lm:phi1phi3} because in all error terms we can estimate the factor $|\nu (x-y)| \lesssim \| \nu \|_\infty \lesssim \rho$ appearing in $\phi_2$, since $Vf_\ell^2$ can be used to integrate in the $x-y$ variable.  
\end{proof} 

\begin{lemma} We have 
\[ \frac{1}{L^3} \text{Re } \int dx dy \, V(x-y) \langle \phi_2 (x;y) , \phi_5 (x;y) \rangle \leq C \rho^{5/2+\eps}  \]
if $\eps, \delta > 0$ are small enough. 
\end{lemma} 
\begin{proof} 
From the definition of $\phi_2, \phi_5$, we find 
\[ \begin{split} 
 \int dx dy \, &V(x-y) \langle \phi_2 (x;y) , \phi_5 (x;y) \rangle \\ = \; & \int dx dy dz dw \, V(x-y) f_\ell^2 (x-y) \nu (x-y) \nu (x-z) \nu (y-w) f_\ell^2 (x-z) f_\ell^2 (y-z) 
 \\  &\hspace{2cm} \times f_\ell^2 (x-w) f_\ell^2 (y-w) f_\ell(z-w) \langle a_z a_w \Psi, J^2 (x) J^2 (y) J(z) J(w) \Psi \rangle. \end{split} \]
We proceed as in the proof of Lemma \ref{lm:phi1phi5} to get rid of the $J$-operators and of the $f_\ell$-factors. After estimating $|\nu (x-y)| \leq \| \nu \|_\infty \lesssim \rho$, all these error terms can be handled exactly as those arising in Lemma \ref{lm:phi1phi5}. Writing $a_z = \sqrt{\rho_0} + b_z$, $a_w = \sqrt{\rho_0} + b_w$ and observing that the contribution of the factors $\sqrt{\rho_0}$ vanishes, because $\int \nu = 0$, we obtain 
\[ \begin{split}  \frac{1}{L^3} &\text{Re} \int dx dy \, V(x-y) \langle \phi_2 (x;y) , \phi_5 (x;y) \rangle \\  \leq\;& C \rho^{5/2+\eps} + \frac{1}{L^3} \int dx dy dz dw \, V(x-y) f_\ell^2 (x-y) \nu (x-y) \nu (x-z) \nu (y-w) \langle b_z b_w \Psi , \Psi \rangle \end{split} \]
if $\eps , \delta > 0$ are small enough. Next we write $b_z = c (\gamma_z) - c^* (\sigma_z)$, $b_w = c(\gamma_w) - c^* (\sigma_w)$. The contribution of all normal ordered terms can bounded analogously as we did with the last term on the r.h.s. of (\ref{eq:IVa}) (again, after estimating $|\nu (x-y)| \lesssim \rho$). We conclude that 
\[  \begin{split}  \frac{1}{L^3} \text{Re } &\int dx dy \, V(x-y) \langle \phi_2 (x;y) , \phi_5 (x;y) \rangle \\ \leq\;& C \rho^{5/2+\eps} + \frac{1}{L^3} \int dx dy dz dw \, V(x-y) f_\ell^2 (x-y) \nu (x-y) \nu (x-z) \nu (y-w) \langle \gamma_z ; \sigma_w \rangle \\  \leq\;& C \rho^{5/2+\eps} + \int dx \, V(x) f_\ell^2 (x) \nu (x) (\nu * \gamma *\sigma *\nu) (x) \\ \leq\;& C \rho^{5/2+\eps} + \int dx \, V(x) f_\ell^2 (x) \nu (x) (\nu * \sigma * \sigma) (x) \end{split} \]
if $\eps , \delta > 0$ are small enough. As argued in the proof of Lemma \ref{lm:phi2}, we have 
\[ |\nu (x)| = |(\gamma^{-1} * \sigma ) (x) | \leq \| (\gamma^{-1} - 1)*\sigma \|_\infty + |\sigma (x)| \lesssim \rho^{3/2} \] 
on the support of $V$. Since $\nu * \sigma * \sigma = \gamma * \sigma - \nu$, we similarly find 
\[ | (\nu * \sigma * \sigma) (x)| \lesssim \rho^{3/2} \]
on the support of $V$. We conclude that 
\[ \frac{1}{L^3} \text{Re} \int dx dy \, V(x-y) \langle \phi_2 (x;y) , \phi_5 (x;y) \rangle \leq C \rho^{5/2+\eps} \]
if $\eps , \delta > 0$ are small enough. 
\end{proof} 

\begin{lemma} We have 
\[ \frac{1}{L^3} \text{Re } \int dx dy \, V(x-y) \langle \phi_3 (x;y) , \phi_4 (x;y) \rangle \leq \rho_0 \int dx \, V(x) f_\ell^2 (x) (\sigma* \sigma )( x) + C \rho^{5/2+\eps} \] if $\eps , \delta > 0$ are small enough. 
\end{lemma} 
\begin{proof} 
From the definition of $\phi_3, \phi_4$, we find 
\[ \begin{split}  \int \; &dx dy \, V(x-y) \langle \phi_3 (x;y) , \phi_4 (x;y) \rangle \\ =\;& \rho_0 \int dx dy dz dw \, V(x-y) f_\ell^2 (x-y) \nu (y-z) \nu (x-w) \\ &\hspace{1cm} \times  f_\ell(x-z) f_\ell(y-z) f_\ell(x-w) f_\ell(y-w) \langle J(x) J(y) J(z) \Psi , a_z a_w^* J(x) J(y) J(w) \Psi \rangle \,. \end{split} \]
With $a_z a_w^* = a_w^* a_z + \delta (z-w)$, we obtain 
\[  \int dx dy \, V(x-y) \langle \phi_3 (x;y) , \phi_4 (x;y) \rangle = \text{D}_1 + \text{D}_2 \]
with 
\[ \begin{split} \text{D}_1 =\;& \rho_0 \int dx dy dz \, V(x-y) f_\ell^2 (x-y) \nu (x-z) \nu (y-z) f_\ell^2(x-z) f_\ell^2(y-z)\\
&\hspace{2cm}\times  \langle \Psi, J(x)^2 J(y)^2 J(z)^2 \Psi \rangle  \\ 
\text{D}_2 =\;& \rho_0 \int dx dy dz dw \, V(x-y) f_\ell^2 (x-y) \nu (x-w) \nu (y-z) f_\ell^2 (x-z) f_\ell^2 (x-w) f_\ell^2 (y-z) \\
&\hspace{2.6cm} \times f_\ell^2 (y-w) f_\ell^2 (z-w) \langle a_w \Psi, J(x)^2 J(y)^2 J(z) J(w) a_z \Psi \rangle . \end{split} \]
To estimate $\text{D}_1$, we observe, with (\ref{eq:zeta-combi1}) and Lemma \ref{lm:a-bds}, that 
\[ \begin{split} 
\Big| \rho_0 &\int dx dy dz \, V(x-y) f_\ell^2 (x-y) \nu (x-z) \nu (y-z) \\
&\hspace{1cm}\times f_\ell^2(x-z) f_\ell^2(y-z)  \langle \Psi, \big( 1 - J^2 (x) J^2 (y) J^2 (z) \big) \Psi \rangle \Big| \\ \leq \; &
\rho \int dx dy dz ds \, V(x-y)f_\ell^2(x-y) |\nu (x-z)|^2 \big (u_\ell (x-s) + u_\ell (y-s) + u_\ell (z-s) \big) \| a_s \Psi \|^2 \\ \lesssim \; &\rho \| \nu \|_2^2 \| u_\ell \|_1 \|Vf_\ell^2\|_1\int ds \, \| a_s \Psi \|^2 \lesssim \rho^{7/2-3\eps-2\delta} L^3 \lesssim \rho^{5/2+\eps} L^3 \end{split} \]
if $\eps , \delta > 0$ are small enough. Moreover,
\begin{equation*}
    \begin{split}
        \bigg|\rho_0\int dxdydz\,& V(x-y) f_\ell^2(x-y) \nu(x-z)\nu(y-z)\big[ 1-f_\ell^2(x-z) f_\ell^2(y-z) \big]\bigg|\\
        \lesssim\;& \rho_0 \|\nu\|_\infty^2 \int dxdydz\,V(x-y) f_\ell^2(x-y) \big[ u_\ell (x-z) + u_\ell (y-z)\big] \lesssim \rho^{5/2+\varepsilon} L^3.
    \end{split}
\end{equation*}
Hence
\begin{equation} \label{eq:D1-bd} \text{D}_1 \leq \rho_0 L^3 \int dx \, V(x) f_\ell^2 (x) (\nu * \nu) (x) + C \rho^{5/2+\eps} L^3 \, .\end{equation} 
As for $\text{D}_2$, we notice that 
 \[  \begin{split} \Big| \rho_0 \int &dx dy dz dw \, V(x-y) f_\ell^2 (x-y) \nu (x-w) \nu (y-z) f_\ell^2 (x-z) f_\ell^2 (x-w) f_\ell^2 (y-z) \\ &\hspace{.5cm} \times f_\ell^2 (y-w) f_\ell^2 (z-w) \langle a_w \Psi, \big( 1 - J(x)^2 J(y)^2 J(z) J(w) \big) a_z \Psi \rangle \Big| \\ \leq \; &C \rho^{5/2+\eps} L^3 + C \rho \int_{|y-z| \leq R}  dx dy dz dw ds \, V(x-y)f_\ell^2(x-y) |\nu (x-w)|^2  \\ &\hspace{3cm} \times \big( u_\ell (x-s) + u_\ell (y-s) + u_\ell (z-s) + u_\ell (w-s) \big) \| a_s a_z \Psi \|^2 
 %\\ \leq \; &C \rho^{5/2+\eps} L^3 + C \rho \| \nu \|_2^2  \int_{|s-z| \leq R} ds dz \, \| a_s a_z \Psi \|^2 + C \rho R^3 \| \nu \|_2^2 \int_{|s-z| \leq \ell} ds dz \, \| a_s a_z \Psi \|^2 
 \\ \leq \; &C \rho^{5/2+\eps} L^3 
  \end{split} \]
  and that 
  \[ \begin{split} \Big| \rho_0 \int &dx dy dz dw \, V(x-y) f_\ell^2 (x-y) \nu (x-w) \nu (y-z)  \\ &\hspace{.5cm} \times \big( 1 - f_\ell^2 (x-z) f_\ell^2 (x-w) f_\ell^2 (y-z) f_\ell^2 (y-w) f_\ell^2 (z-w) \big)  \langle a_w \Psi, a_z \Psi \rangle \Big| \\ \leq \; &C \rho^{5/2+\eps} L^3 + C \rho \int  dx dy dz dw \, V(x-y) f_\ell^2(x-y)|\nu (x-w)| |\nu (y-z)|   \\ &\hspace{.5cm} \times \big( u_\ell (x-z) + u_\ell (x-w) + u_\ell (y-z) + u_\ell (y-w) + u_\ell (z-w)  \big) \| a_z \Psi \|^2 
 %\\ \leq \; &C \rho^{5/2+\eps} L^3 + C \rho \| \nu \|_2^2  \int_{|s-z| \leq R} ds dz \, \| a_s a_z \Psi \|^2 + C \rho R^3 \| \nu \|_2^2 \int_{|s-z| \leq \ell} ds dz \, \| a_s a_z \Psi \|^2 
 \\ \leq \; &C \rho^{5/2+\eps} L^3 
  \end{split} \]
if $\eps , \delta > 0$ are small enough. Hence
\[ \begin{split}  \text{D}_2 \leq \; &\rho_0 \int dx dy dz dw \, V(x-y) f_\ell^2 (x-y) \nu (x-w) \nu (y-z) \langle b_w \Psi, b_z \Psi \rangle + C \rho^{5/2+\eps} \\ 
\leq \; &\rho_0 \int dx dy dz dw \, V(x-y) f_\ell^2 (x-y) \nu (x-w) \nu (y-z) \langle \sigma_w ; \sigma_z \rangle \\ &+ 
\rho_0 \int dx dy dz dw \, V(x-y) f_\ell^2 (x-y) \nu (x-w) \nu (y-z) \\ &\hspace{1cm} \times \langle \Psi, \big( c^* (\gamma_w) c (\gamma_z) + c^* (\sigma_z) c (\sigma_w) + c^* (\gamma_w) c^* (\sigma_z) + c (\sigma_w) c (\gamma_z) \big) \Psi \rangle \\ &+C \rho^{5/2+\eps} .
\end{split} \] 
With (\ref{eq:zeta-combi1}) and Lemma \ref{lm:c's}, we find 
\[ \rho \int dx dy dz dw \, V(x-y)f_\ell^2(x-y) |\nu (x-w)| |\nu (y-z)| \big[ \| c (\gamma_z) \Psi \|^2 + \| c (\sigma_z) \Psi \|^2 \big] \leq C \rho^{5/2+\eps} L^3 \]
and 
\[ \begin{split} \rho \int dx dy dz &dw \, V(x-y) f_\ell^2(x-y) |\nu (x-w)| |\nu (y-z)|  \| c (\gamma_w) c (\sigma_z) \Psi \| \\ \leq \; &\rho \Big( \int dx dy dz dw \, V(x-y)f_\ell^2(x-y) |\nu (x-w)| \|c (\gamma_w) c (\sigma_z) \Psi \|^2 \Big)^{1/2} \\ &\hspace{2cm} \times  \Big(  \int dx dy dz dw \, V(x-y) f_\ell^2(x-y)|\nu (x-w)|  |\nu (y-z)|^2 \Big)^{1/2} \\ \leq \; &C \rho^{5/2+\eps} L^3 \end{split} \]
if $\eps, \delta > 0$ are small enough. Therefore
\[ \begin{split}  \text{D}_2 \leq \; &\rho_0 \int dx\, V(x) f_\ell^2 (x)  (\nu * \nu * \sigma * \sigma) (x) + C \rho^{5/2+\eps} L^3. \end{split} \] 
With (\ref{eq:D1-bd}) and since $1 + \sigma*\sigma = \gamma * \gamma$, we conclude that 
\[\frac{1}{L^3}  \int dx dy \, V(x-y) \langle \phi_3 (x,y) \Psi ; \phi_4 (x,y) \Psi \leq \rho_0  \int dx \, V(x) f_\ell^2 (x) (\sigma * \sigma) (x) + C \rho^{5/2+\eps}.\] 
\end{proof} 

\begin{lemma}\label{lm:phi3phi5} For $j=3,4$, we have 
\[ \frac{1}{L^3} \text{Re } \int dx dy \, V(x-y) \langle \phi_j (x;y) , \phi_5 (x;y) \rangle \leq C \rho^{5/2+\eps/2} \]
if $\eps , \delta > 0$ are small enough. 
\end{lemma} 
\begin{proof} 
From Lemma \ref{lm:phi3}, we have 
\[ \frac{1}{L^3} \int dx dy \, V(x-y) \| \phi_j (x;y) \|^2 \leq C \rho^{5/2} \]
for $j=3,4$. From Lemma \ref{lm:phi5}, we also find 
\[ \frac{1}{L^3} \int dx dy \, V(x-y) \| \phi_5 (x;y) \|^2 \leq C \rho^{5/2+\eps} \]
for $\eps, \delta > 0$ small enough. By Cauchy-Schwarz, we conclude that 
\[ \frac{1}{L^3} \Big| \int dx dy \, V(x-y) \langle \phi_j (x,y) ; \phi_5 (x,y) \rangle \Big| \leq C \rho^{5/2+\eps/2}  \]
for $\eps , \delta > 0$ small enough. 
\end{proof} 

Combining Lemmas \ref{lm:phi1} - \ref{lm:phi3phi5}, we obtain the desired upper bound on the potential energy in the trial state $\Psi$.
\begin{proof}[Proof of Prop. \ref{prop:pot}] 
Collecting all estimates from Lemmas \ref{lm:phi1} - \ref{lm:phi3phi5}, we arrive at 
\[ \begin{split} 
\frac{1}{2L^3} \int &dx dy \, V(x-y) \| a_x a_y \Psi \|^2 \\ \leq \; &\frac{\rho_0^2}{2} \int dx \, V(x) f_\ell^2 (x) + \rho_0 \| \sigma \|_2^2 \int dx \, V(x) f_\ell^2 (x) \\ &+ \rho_0 \int dx \, V(x) f_\ell^2 (x) \nu (x) + \rho_0 \int dx V(x) f_\ell^2 (x) (\nu * \sigma * \sigma) (x) \\ &+ \rho_0 \int dx V(x) f_\ell^2 (x) (\sigma * \sigma) (x) + C \rho^{5/2+ \eps/2} .
\end{split} \]
With $\nu + \nu * \sigma * \sigma = \sigma * \gamma$, we obtain 
\[ \begin{split} 
\frac{1}{2L^3} \int &dx dy \, V(x-y) \| a_x a_y \Psi \|^2 \\ \leq \; &\frac{\rho_0^2}{2} \int dx \, V(x) f_\ell^2 (x) + \rho_0 \| \sigma \|_2^2 \int dx \, V(x) f_\ell^2 (x) \\ &+ \rho_0 \int dx \, V(x) f_\ell^2 (x)  \big( (\sigma * \sigma) (x) + (\sigma * \gamma) (x) \big) + C \rho^{5/2+ \eps/2} .
\end{split} \]
In the last term, we decompose $\sigma * \gamma = \sigma * (\gamma-1) + \sigma$ and we observe that, on the support of $V$, 
\[ |\sigma (x)| \leq \ell_0^{-3} |\ph (x/\ell_0)| \Big| \int \tilde{\sigma} (z) dz \Big| \lesssim \rho^{3/2+3\eps/2}. \]
Hence 
\[ \begin{split} 
\frac{1}{2L^3} \int &dx dy \, V(x-y) \| a_x a_y \Psi \|^2 \\ \leq \; &\frac{\rho_0^2}{2} \int dx \, V(x) f_\ell^2 (x) + \rho_0 \| \sigma \|_2^2 \int dx \, V(x) f_\ell^2 (x) \\ &+ \rho_0 \int dx \, V(x) f_\ell^2 (x)  \big( (\sigma * \sigma) (x) + (\sigma * (\gamma-1)) (x) \big) + C \rho^{5/2+ \eps/2} .
\end{split} \]

%{\bf [The expression above is more convenient, if we want to %compare with \cite[Proposition 2.5]{BBCOS}, see Proposition 
%\ref{prop:Psi-energy}]}
%
%Switching to Fourier space and estimating $|e^{ip\cdot x} - %1| \leq C |p| |x|$, we find   
%\[ \big| (\sigma * \sigma) (x) - (\sigma * \sigma)(0) \big| 
%\leq C |x| \frac{1}{|\Lambda|} \sum_{p \in \Lambda^*} |p| |
%\hat{\sigma}_p|^2 \leq C |x| \| \sigma \|_2 \| \nabla \sigma 
%\|_2 \lesssim |x| \rho^{7/4} \]
%and, similarly, 
%\[  \big| (\sigma * (\gamma-1)) (x) - (\sigma * (\gamma-1))
%(0) \big| \leq C |x| \rho^{7/4}. \]
%We conclude that  
%\[ \begin{split} 
%\frac{1}{2L^3} &\int dx dy \, V(x-y) \| a_x a_y \Psi \|^2 \\ 
%\leq \; &\frac{\rho_0^2}{2} \int dx \, V(x) f^2 (x) +  %2\rho_0 \| \sigma \|_2^2 \int dx \, V(x) f^2 (x)  + \rho_0 %(\sigma * (\gamma-1)) (0) \int dx \, V(x) f^2 (x) \\ &+ C 
%\rho^{5/2+ \eps/2} .
%\end{split} \]
\end{proof}

%From the definition of $\phi_3, \phi_5$, we find 
%\[ \begin{split} \int dx dy \, &V(x-y) \langle \phi_3 (x,y) ; \phi_5 (x,y) \rangle \\ &= \sqrt{\rho_0} \int dx dy dz ds dw \, V(x-y) f^2 (x-y) \nu (y-z) \nu (x-s) %J(x) J(y) J(s) J(w) \Psi \rangle \end{split} \]
%From the canonical commutation relations, we have 
%\[ a_z a_s^* a_w^* = \delta (z-s) a_w^* + \delta (z-w) a_s^* + a_s^* a_w^* a_z \]
%\[ \begin{split} \int dx &dy \, V(x-y) \langle \phi_3 (x,y) ; \phi_5 (x,y) \rangle \\ = \; &\sqrt{\rho_0} \int dx dy dz dw \, V(x-y) f^2 (x-y) \nu (y-z) \nu (x-z) 
%\nu (y-w) \\ &\hspace{.2cm} \times f^2 (x-z) f^2 (y-z) f^2 (x-w) f^2 (y-w) f^2 (z-w) \langle a_w \Psi, J^2 (x) J^2 (y) J^2 (z) J (w) \Psi \rangle \\ &+  
%\sqrt{\rho_0} \int dx dy dz ds \, V(x-y) f^2 (x-y) \nu (x-s) \nu^2 (y-z) \\ &\hspace{.2cm} \times f^2 (x-z) f^2 (y-z) f^2 (x-s) f^2 (y-s) f^2 (z-s) \langle a_s %\Psi, J^2 (x) J^2 (y) J^2 (z) J (s) \Psi \rangle \\ &+ \sqrt{\rho_0} \int dx dy dz ds dw \, V(x-y) f^2 (x-y) \nu (y-z) \nu (x-s) \nu (y-w) \\ &\hspace{.2cm} 
%\times  f^2 (y-s) f^2 (x-w) f^2 (y-w) f^2 (w-z) f^2 (s-z) f (s-w) \\ &\hspace{6cm} \times  \langle a_s a_w \Psi, J^2 (x) J^2 (y) J(z) J(s) J(w) a_z \Psi 
%\rangle \\ =: \; &\text{E}_1 + \text{E}_2 + \text{E}_3  \end{split} \]
%To bound $\text{E}_1$, we observe that 
%\[  \begin{split} \Big| \sqrt{\rho_0} &\int dx dy dz dw \, V(x-y) f^2 (x-y)  \nu (y-z) \nu (x-z) \nu (y-w)  f^2 (x-z) f^2 (y-z) \\ &\hspace{.2cm} \times f^2 (x-w) %f^2 (y-w) f^2 (z-w) \langle a_w \Psi,  \big(1- J^2 (x) J^2 (y) J^2 (z) J (w) \big) \Psi \rangle \Big| \\  \leq \; &\sqrt{\rho_0} \Big( \int dx dy dz dw \, V(x-y) |%\nu (y-z)| \\ &\hspace{2cm} \times  \big( u (s-x) + u (s-y) + u(s-z) + u (s-w) \big) \| a_s a_w \Psi \|^2 \Big)^{1/2} \\ &\times \Big( \int dx dy dz dw \, V(x-%y) |\nu (y-z)|  |\nu (x-z)|^2 |\nu (y-w)|^2 \\ &\hspace{2cm} \times \big( u (s-x) + u (s-y) + u(s-z) + u (s-w) \big) \| a_s \Psi \|^2 \Big)^{1/2} \\ \leq \; &C 
%\rho^{5/2+\eps}  \end{split} \]
%if $\eps , \delta > 0$ are small enough. 
%\end{proof} 


\appendix

\section{Properties of the kernels: proof of Lemma \ref{lm:eta}}
\label{app:kernels}

We proceed similarly as in the proof of \cite[Lemma 2.2]{BBCOS}. The main difference is that now we can only rely on the estimates $\| \nabla^m \widehat{V}_\text{eff} \|_\infty \leq C_m$, for every $m \in \bN$, from Lemma \ref{Lem:general_potential_eff} (in \cite{BBCOS}, thanks to the explicit form of the effective potential, we could use the stronger bound $|\nabla^m \widehat{V}_\text{eff} (k)| \lesssim |k|^{-1}$, for $|k| \geq 1$). 

Following the strategy of \cite{BBCOS}, in the first part of the proof we bound the Fourier coefficients $\widehat{s}_k$ and their (discrete) derivatives $\delta_j^m \widehat{s}_k$. We find (proceeding as in the proof of \cite[Lemma 2.2]{BBCOS}, recalling that $\widehat{V}_\text{eff} (0) = 8 \pi \frak{a} > 0$ and therefore that $\widehat{V}_\text{eff} (k) \geq c > 0$, for $k$ small enough) 
%\begin{equation} \label{eq:sfourier} |\widehat{s}_k| \leq C \min \Big\{ % \frac{\rho^{1/4}}{|k|^{1/2}} , \frac{\rho}{k^2} \Big\} \end{equation} 
%for all $k \in \Lambda^*_+$, as well as
\begin{align}
\label{eq:sfourier_I}
    |\delta_j^m \widehat{s}_k| \leq C_m \, |k|^{-m} \min \Big\{  \frac{\rho^{1/4}}{|k|^{1/2}} , \frac{\rho}{k^2} \Big\} 
\end{align}
for all $m\in \mathbb N$ and $k\in \Lambda^*_+$ with $|k|\leq 1$. For $|k| \geq 1$, we get $|\delta_j^m \widehat{s}_k| \leq C_m \rho / k^2$ for all $m \in \bN$ (this should be compared with \cite[(A.8)]{BBCOS}, where we could prove the stronger estimate $|\delta_j^m \widehat{s}_k| \leq C_m \rho / |k|^3$, using the decay of $\widehat{V}_\text{eff}$). For $|k| \geq 1$, it is also important to derive the improved estimate 
\begin{align}
\label{eq:sfourier_II}
   \left| \delta_j^m\left[ \widehat{s}_k+\frac{\rho_0 \widehat{V}_\mathrm{eff}(k)}{2|k|^2}-\frac{\rho_0^2 \widehat{V}_\mathrm{eff}(k)^2}{2|k|^4}+\frac{11 \rho_0^3 \widehat{V}_\mathrm{eff}(k)^3}{16|k|^6} \right]\right| & \leq C_m \frac{\rho^4}{|k|^8} 
\end{align}
for all $m\in \bN$. Compared with \cite[(A.9)]{BBCOS}, here it is important to expand $\widehat{s}_k$ to higher order, to compensate the lack of decay of $\widehat{V}_\text{eff}$. 

Next, we use the bounds for the Fourier coefficients $\widehat{s}_k$ to estimate $s$ in position space. In particular, we show that 
\begin{equation}\label{eq:s-point} | s (x)| \leq C \frac{\rho}{|x|} \min \Big\{ 1 , \frac{1}{(\rho^{1/2} |x|)^{3/2}} \Big\} \,, \qquad |\nabla s (x)| \leq C \frac{\rho |\hspace{-.05cm} \log \rho |}{x^2} \end{equation} 
for all $|x| \geq  15 r_0$, and that, for $r\geq 15r_0$
\begin{equation}\label{eq:int-s}\begin{split} 
    \int_{|x|\geq r}  |\nabla s(x)|^2dx  &\lesssim \min\left\{r^{-4}\rho^{\frac{1}{2}},r^{-1}\rho^2\right\},\\
       \int_{|x|\geq r}   \left|\Delta s (x) \right|^2dx& \lesssim   r^{-3}\rho^2,\\
     \int_{|x|\geq r}  \left|\nabla \Delta s (x) \right|^2dx& \lesssim   r^{-5}\rho^2.
     \end{split} 
\end{equation}
%as well as the bounds
%\begin{equation}\label{eq:est-combi}
% \int_{|x|\geq \ell_0} |s(x)|^2 dx  \leq C\ell_0^{-2}\rho^{\frac{1}{2}} % , \;   \int_{15r_0 \leq |x|\leq  2\ell} 
% |s(x)|^2 dx  \leq C\ell\rho^{2} . \end{equation} 

To show (\ref{eq:s-point}) and (\ref{eq:int-s}), we decompose $s$ into a low- and a high-momentum component. To this end, we choose $\chi \in C_c^\infty (\bR^3)$ with $\chi (z) = 1$ for $|z| \leq 2$ and $\chi (z) = 0$, for $|z| \geq 4$. Moreover, for $t > 0$, we set $\chi_t (z) = \chi (z/t)$ and $s_t = \widecheck{\chi}_t * s$. With this notation, we write $s = s_1 + (s-s_1)$. The contribution of the low-momentum part $s_1$ to (\ref{eq:s-point}), (\ref{eq:int-s}) can be controlled exactly as in \cite[Lemma 2.2]{BBCOS} (at low-momenta, the analysis relies on (\ref{eq:sfourier_I})); we skip the details. To estimate the high momentum component $s-s_1$, on the other hand, we have to modify the approach of \cite{BBCOS}. We introduce $h = - \rho_0 \chi_{r_0} / 8\pi |.| \in L^2 (\bR^3)$ and we define $D' \in L^2 (\bR^3)$ through its Fourier transform 
\begin{align*}
    \widehat{D}'(k):=\widehat{V}_\mathrm{eff}(k)\widehat{h}(k).
\end{align*}
%For a mollifier $m_\epsilon(y):=\epsilon^{-3}m(\epsilon^{-1} y)$, where %$m$ is smooth and compactly supported with $\int m(y) \mathrm{d}y=1$, 
%we clearly have $m_\epsilon * D'\underset{\epsilon \rightarrow 0}
%{\longrightarrow }D'$ in $L^2(\mathbb R^3)$, and as a consequence there %exists a subsequence $\epsilon_j$ that is almost everywhere convergent %to $D'$. 
Recalling that $V_\mathrm{eff}$ is supported in $B_{r_0}(0)$, 
%see Lemma \ref{Lem:general_potential_eff}, an explicit computation %demonstrates for $|x|>r_0$
%\begin{align*}
%   m_\epsilon * D'(x)=-\frac{\rho_0}{8\pi}\int  \frac{\chi_{r_0}(x-y)}{|%x-y|}(m_\epsilon * V_\mathrm{eff})(y)dy\underset{\epsilon \rightarrow 0}
%{\longrightarrow } -\frac{\rho_0}{8\pi}\int  \frac{\chi_{r_0}(x-y)}{|x-%y|} V_\mathrm{eff}(y)dy,
%\end{align*}
we find, for (almost every) $|x|>r_0$, 
\begin{align}
\label{Eq:explicit_formula_D(x)}
    D'(x) = -\frac{\rho_0}{8\pi}\int  \frac{\chi_{r_0}(x-y)}{|x-y|}V_\mathrm{eff}(y)dy.
\end{align}
In Fourier space, at high momenta, $D'$ is an approximation for the factor $- \rho_0 \widehat{V}_\text{eff} (k) / 2k^2$ appearing in the expansion of $\widehat{s}_k$ (see (\ref{eq:sfourier_II})). More precisely, we find 
\begin{align*}
    \widehat{D}'_k +\frac{\rho_0 \widehat{V}_\mathrm{eff}(k)}{2|k|^2}=\lim_{R\rightarrow \infty} \frac{\rho_0 \widehat{V}_\mathrm{eff}(k)}{8\pi}\int \frac{\chi_R(x)-\chi_{r_0}(x)}{|x|} \, e^{i k \cdot x}dx.
\end{align*}
From the smoothness of $(\chi_R(x) - \chi_{r_0}(x))/|x|$, in particular from the estimate \[ \int \left|\nabla^\alpha \left(\frac{\chi_R(x)-\chi_{r_0}(x)}{|x|}\right)\right| dx\leq C \] holding for all $\alpha  \geq 3$, uniformly in $R>0$, we obtain, for any $\alpha,m\in \mathbb N$, 
\begin{equation}\label{eq:D'+}
\left|\delta^m_j \left[\widehat{D}'_k +\frac{\rho_0 \widehat{V}_\mathrm{eff}(k)}{2|k|^2}\right]\right|
%=\left|\delta^\alpha_j \left[\widehat{V}_\mathrm{eff}(k)\widehat{h}(k) +%\frac{\rho_0 \widehat{V}_\mathrm{eff}(k)}{2|k|^2}\right]\right|
\lesssim \rho |k|^{-\alpha}
\end{equation}
for every $|k| \geq 1$. From \eqref{eq:sfourier_II}, we conclude that 
\[ \Big| \delta_j^m \Big[ \widehat{s}_k - \widehat{D}'_k - 2 \widehat{D}^{'2}_k - \frac{11}{2} \widehat{D}^{'3}_k \Big] \Big| \lesssim \frac{\rho}{|k|^8} \,. \]
To approximate the high momentum part $s-s_1$, we set $D := D' - \widecheck{\chi}_1 * D'$ (in momentum space $\widehat{D}_k = (1- \chi_1 (k)) \widehat{D}'_k$). We define the error term $E$, requiring that 
\begin{align}
\label{Eq:decomposition_s_high_momenta}
   s-s_1 = D + 2 D*D + \frac{11}{2}D*D*D + E \,.
\end{align}
Observing that $\widehat{E}_k = \widehat{s}_k - \widehat{D}'_k - 2 \widehat{D}^{'2}_k - 11 \widehat{D}^{'3}_k /2$ for $|k| \geq 4$ (because $\widehat{D}_k = \widehat{D}'_k$, for $|k| \geq 4$) and using (\ref{eq:D'+}) and $|\delta_j^m \widehat{s}_k| \lesssim \rho / k^2$ for $1 \leq |k| \leq 4$, we conclude that  
\begin{equation}\label{eq:deltaE} |\delta_j^m \widehat{E}_k| \lesssim \frac{\rho}{|k|^8} \end{equation}
for all $|k| \geq 1$. Furthermore, $\widehat{E}_k = 0$, for $|k| \leq 1$. With the estimate 
 \begin{equation*}
        |x_j^m f(x)|\le \Big(\frac{\pi}{2}\Big)^m\Big| \Big(\frac{L}{2\pi}\Big)^m\big(e^{-i\frac{2\pi}{L}x_j}-1\big)^m f(x)\Big|=\Big(\frac{\pi}{2}\Big)^m  \big| (\widecheck{\delta^m_j \widehat{f}} ) (x)\big|\leq \Big(\frac{\pi}{2}\Big)^m \left\|\delta^m_j \widehat{f} \right\|_1 
    \end{equation*}
we can translate (\ref{eq:deltaE}) to position space. We find the pointwise bounds
\[ |E(x)|, |\nabla E(x)| , |\Delta E (x)| , |\nabla \Delta E (x)| \lesssim \frac{\rho}{|x|^m}\]
for any $m \geq 1$ (the bound for $\nabla \Delta E$ requires $k \to \delta_j^m |k|^3 \widehat{E} (k)$ to be summable; that's why we had to expand $s$ to third order). To estimate $s-s_1$, we still have to control the contributions $D, D*D, D*D*D$ appearing on the r.h.s. of  \eqref{Eq:decomposition_s_high_momenta}. Since $V_\mathrm{eff}$ is supported in $B_{r_0}(0)$ and since the function $\chi_{r_0}$ appearing in \eqref{Eq:explicit_formula_D(x)} is supported in $B_{4r_0}(0)$, we conclude that $D'$ is supported in $B_{5r_0}(0)$. Hence, for $|x| \geq 5r_0$ and $q\in \mathbb N$, we find that
    \begin{align*}
       | \nabla^{q} D(x) | = | -\left(\nabla^q \widecheck{\chi}_{1}\right)*D' | = \Big| \frac{\rho_0}{8\pi}\int  \frac{\chi_{r_0}(x-y)}{|x-y|}(\nabla^q \widecheck{\chi}_{1} * V_\mathrm{eff})(y) dy \Big| \lesssim \frac{\rho}{|x|^m}
    \end{align*}
for any $m \in \bN$. Similarly, we can estimate $\nabla^q (D*D)$ and $\nabla^q (D*D*D)$, for $|x|\geq 10r_0$ and, respectively, for $|x|\geq 15r_0$. By \eqref{Eq:decomposition_s_high_momenta}, in combination with the pointwise bounds for $E$ and its derivatives, we therefore obtain, for $|x|\geq 15r_0$, 
    \begin{align*}
        \left|(s-s_1)(x)\right|, \left|\nabla (s-s_1)(x)\right|, \left|\Delta(s-s_1)(x)\right|, \left|\nabla \Delta (s-s_1)(x)\right|\lesssim \frac{\rho}{|x|^{m}} 
    \end{align*}
for any $m > 0$. Combining these bounds with the estimates for the low-momentum component $s_1$ (established as in \cite{BBCOS}) we obtain (\ref{eq:s-point}), (\ref{eq:int-s}). The rest of the proof is identical to \cite[Lemma 2.2]{BBCOS} (the bounds (\ref{eq:int-s}) play an important role in the proof of (\ref{eq:decay})). 


\begin{thebibliography}{55}

\def\bibskip{\\[-0.6cm]}

\bibitem{BBCOS} G. Basti, M. Brooks, S. Cenatiempo, A. Olgiati, B. Schlein. The Lee-Huang-Yang energy for a dilute gas of hard spheres: an upper bound. Preprint arXiv:2603.13084. \bibskip


\bibitem{BCGOPS} G. Basti, S. Cenatiempo, A. Giuliani, A. Olgiati, G. Pasqualetti, B. Schlein. Upper bound for the ground state energy of a dilute Bose gas of hard spheres.
{\it Arch. Ration. Mech. Anal.} {\bf 248} (2024), no. 100. \bibskip

\bibitem{BCOPS} G. Basti, S. Cenatiempo, A. Olgiati, G. Pasqualetti, B. Schlein. A Second Order Upper Bound for the Ground State Energy of a Hard-Sphere Gas in the Gross–Pitaevskii Regime. {\it Commun. Math. Phys.} \textbf{399} (2023), 1–55. \bibskip

\bibitem{BCS}
G. Basti. S. Cenatiempo, B. Schlein. A new second order upper bound for the ground state energy of dilute Bose gases. {\it Forum Math. Sigma} {\bf 9} (2021), no. e74. \bibskip

\bibitem{BBCS1}
C. Boccato, C. Brennecke, S. Cenatiempo, B. Schlein. Optimal rate for Bose-Einstein condensation in the Gross-Pitaevskii regime. {\it Commun. Math. Phys} {\bf 376} (2020), 1311--1395. \bibskip

\bibitem{BBCS2}
C. Boccato, C. Brennecke, S. Cenatiempo, B. Schlein. Bogoliubov Theory in the Gross-Pitaevskii limit. {\it Acta Mathematica} {\bf 222} (2019), no. 2, 219--335. \bibskip

\bibitem{Bog}
N. N. Bogoliubov. On the theory of superfluidity.
{\it Izv. Akad. Nauk. USSR} {\bf 11} (1947), 77. Engl. Transl. {\it J. Phys. (USSR)} {\bf 11} (1947), 23. \bskip

\bibitem{B} M. Brooks. Diagonalizing Bose Gases in the Gross-Pitaevskii Regime and  beyond. {\it Comm. Math. Phys.} {\bf 406} (2025), no. 17. \bibskip

\bibitem{BOSS} 
M. Brooks, J. Oldenburg, D. Saint Aubin, B. Schlein. Third Order Upper Bound for the Ground State Energy of the Dilute Bose Gas. {\it Comm. Math. Phys.} {\bf 407} (2026), no. 163. \bibskip

\bibitem{COSS}
C. Caraci, A. Olgiati, D. Saint Aubin, B. Schlein. Third order corrections to the ground state energy of a Bose gas in the Gross-Pitaevskii regime. {\it Comm. Math. Phys.}  {\bf 406} (2025), no. 153.  \bibskip

\bibitem{Dy}
F.J. Dyson. Ground-State Energy of a Hard-Sphere Gas. {\it Phys. Rev.} {\bf 106} (1957), 20--26. \bibskip

\bibitem{ESY}
L. Erd\H os, B. Schlein, H.-T. Yau. Ground-state energy of a low-density Bose gas: a second order upper bound. {\it Phys. Rev. A} {\bf 78} (2008), 053627. \bibskip

\bibitem{FJGMOT} 
S. Fournais, L. Junge, T. Girardot, L. Morin, M. Olivieri, A. Triay. The free energy of dilute Bose gases at low temperatures interacting via strong potentials. {\it Ann. Henri Poincaré} (2026).\bibskip


\bibitem{FS1}
S. Fournais, J.P. Solovej. The energy of dilute Bose gases. {\it Ann. Math.} {\bf 192} (2020), no. 3, 893--976.  \bibskip

\bibitem{FS2}
S. Fournais, J.P. Solovej. The energy of dilute Bose gases II: The general case.
{\it Invent. Math.} {\bf 232} (2023), 863--994. \bibskip

\bibitem{HHNST} F. Haberberger, C. Hainzl, P. T. Nam, R. Seiringer, A. Triay. The free energy of dilute Bose gases at low temperatures. Preprint arXiv:2304.02405. \bibskip

\bibitem{HHST} F. Haberberger, C. Hainzl, B. Schlein, A. Triay. Upper Bound for the Free Energy of Dilute Bose Gases at Low Temperature. {\it Adv. Math.} {\bf 490} (2026), 110825. \bibskip

\bibitem{HST} C. Hainzl, B. Schlein, A. Triay. Bogoliubov theory in the Gross-Pitaevskii limit: a simplified approach. {\it Forum of Math. Sigma} {\bf 10} (2022), e90. \bibskip


\bibitem{HP}
N. M. Hugenholtz, D. Pines. Ground-State Energy and Excitation Syectrum of a
System of Interacting Bosons. {\it Phys. Rev.} {\bf 116} (1959), no.3,  489–506. 
\bibskip

\bibitem{LHY}
T. D. Lee, K. Huang, and C. N. Yang, Eigenvalues and eigenfunctions of a Bose
system of hard spheres and its low-temperature properties. {\it Phys. Rev.} {\bf 106} (1957), 1135--1145. \bibskip


\bibitem{LY} 
E.~H.~Lieb, J.~Yngvason. Ground State Energy of the low density Bose Gas. {\it Phys. Rev. Lett.} {\bf 80} (1998), 2504--2507.   \bibskip

\bibitem{RT}
L.~Reichmann, A.~Triay. On the spherical symmetry and finite-range assumptions of the interaction potential in the low energy study of dilute Bose gases. Preprint arXiv: 2607.00859. \bibskip


\bibitem{R} D. Ruelle. {\it Statistical Mechanics: Rigorous Results.} World Scientific (1999). \bibskip

\bibitem{Sa}
K. Sawada. Ground-State Energy of Bose-Einstein Gas with Repulsive Interaction. {\it Phys. Rev.} {\bf 116} (1959), no. 6,  1344--1358. \bibskip

\bibitem{Wu}
T.T. Wu. Ground State of a Bose System of Hard Spheres. {\it Phys. Rev.} {\bf 115}, (1959) no.6, 1390--1404.  \bibskip

\bibitem{YY}
H.-T. Yau, J. Yin. The second order upper bound for the ground state energy of a Bose gas. {\it J. Stat. Phys.} {\bf 136} (2009), no. 3, 453--503. \bibskip

\end{thebibliography}
 
\end{document} 



